\chapter{Mathematical Reference}
\label{appA}

\begin{chapterintro}
This appendix collects the mathematics a communications engineer reaches for every week:
Fourier transform properties and pairs in continuous and discrete time, the special
functions that appear in error-rate formulas (the Q-function and its bounds, Marcum's Q,
Bessel and gamma functions), a table of the probability distributions that describe noise
and fading, Gaussian random vectors, random processes and their spectra, the linear algebra
of MIMO and equalisation (including complex gradients), constrained optimisation with
water-filling as the worked example, the decibel and its many dialects, physical constants,
frequency-band names, and the series and integrals that keep recurring. Derivations are
short; the chapters that use each result are cross-referenced so you can see it at work.
Keep a finger in this appendix as you read the book.
\end{chapterintro}

%======================================================================
\section{Notation and conventions}
\label{sec:appA:notation}
%======================================================================

Table~\ref{tab:appA:notation} lists the conventions used throughout the book. Two of them
cause more confusion than all the others combined, so they are worth stating up front.
First, the Fourier transform is written in \emph{hertz}, $X(f)=\int x(t)e^{-j2\pi ft}\,dt$,
not in radians per second. With this choice there are no stray factors of $2\pi$ in the
inverse transform, in Parseval's theorem or in the convolution theorem, and bandwidths read
directly in Hz. Second, $\sinc x=\sin(\pi x)/(\pi x)$ is the \emph{normalised} sinc, whose
zeros fall on the non-zero integers; the unnormalised $\sin x/x$ never appears without being
written out.

\begin{table}[htbp]\centering\small
\caption{Notation used in this book.}\label{tab:appA:notation}
\begin{tabular}{@{}ll@{}}
\toprule
Symbol & Meaning\\
\midrule
$x(t)$, $x[n]$ & continuous-time signal; discrete-time sequence\\
$X(f)$, $X(e^{j\omega})$, $X[k]$, $X(z)$ & Fourier transform; DTFT; DFT; $z$-transform\\
$\tilde x(t)$ or $x_b(t)$ & complex envelope (baseband equivalent): $x(t)=\re\{\tilde x(t)e^{j2\pi f_c t}\}$\\
$*$, $\circledast$ & linear convolution; circular convolution\\
$z^*$, $\vect{A}^T$, $\vect{A}^H$ & complex conjugate; transpose; conjugate (Hermitian) transpose\\
$\vect{a}$, $\vect{A}$, $\vect{I}_n$ & column vector; matrix; $n\times n$ identity\\
$\|\vect{a}\|$, $\|\vect{A}\|_F$ & Euclidean norm; Frobenius norm\\
$\operatorname{tr}$, $\det$, $\operatorname{diag}$, $\operatorname{vec}$, $\otimes$ & trace; determinant; diagonal matrix; column stacking; Kronecker product\\
$\E[\cdot]$, $\operatorname{Var}[\cdot]$, $\Pr\{\cdot\}$ & expectation; variance; probability\\
$\mathcal{N}(\mu,\sigma^2)$, $\mathcal{CN}(\vect{\mu},\vect{C})$ & real Gaussian; circularly symmetric complex Gaussian\\
$\delta(t)$, $\delta[n]$ & Dirac delta; Kronecker delta\\
$u(t)$, $\operatorname{rect}(t)$, $\Lambda(t)$ & unit step; unit rectangle ($|t|<\tfrac12$); unit triangle ($|t|<1$)\\
$\sinc x$ & $\sin(\pi x)/(\pi x)$, with $\sinc 0=1$\\
$\Q(x)$ & Gaussian tail probability $\Pr\{\mathcal N(0,1)>x\}$\\
$\log$, $\ln$, $\log_2$ & base 10 in decibel formulas, natural, base 2 (bits)\\
$E_s$, $E_b$, $N_0$ & energy per symbol; energy per bit; one-sided noise PSD (W/Hz)\\
$T$, $R_s=1/T$, $R_b$ & symbol period; symbol rate (baud); bit rate\\
$(x)^+$ & $\max(x,0)$\\
\bottomrule
\end{tabular}
\end{table}

\paragraph{Complex baseband.} A real bandpass signal centred on $f_c$ is written
$x(t)=\re\{\tilde x(t)e^{j2\pi f_c t}\}=x_I(t)\cos 2\pi f_c t-x_Q(t)\sin 2\pi f_c t$ with
$\tilde x=x_I+jx_Q$. Its spectrum is $X(f)=\tfrac12[\tilde X(f-f_c)+\tilde X^*(-f-f_c)]$, and its
energy is half the energy of the envelope, $\int x^2\,dt=\tfrac12\int|\tilde x|^2\,dt$ (for
$f_c$ well above the bandwidth). Real white noise of two-sided PSD $N_0/2$ becomes complex
baseband noise of PSD $N_0$, with $N_0/2$ in each of $I$ and $Q$
(Chapters~\ref{ch:signals} and~\ref{ch:noise}). Some texts scale the envelope by $\sqrt2$ so
that energies match; this book does not, so watch the factor of two when comparing formulas.

%======================================================================
\section{Fourier analysis}
\label{sec:appA:fourier}
%======================================================================

\subsection{The continuous-time Fourier transform}

The transform pair is
\begin{equation}
X(f)=\int_{-\infty}^{\infty}x(t)\,e^{-j2\pi ft}\,dt,\qquad
x(t)=\int_{-\infty}^{\infty}X(f)\,e^{j2\pi ft}\,df.
\label{eq:appA:ft}
\end{equation}
It exists in the ordinary sense for absolutely integrable or finite-energy signals and, with
Dirac deltas allowed, for periodic signals and constants as well. Tables~\ref{tab:appA:ftprops}
and~\ref{tab:appA:ftpairs} collect the properties and pairs that recur in this book;
Chapter~\ref{ch:signals} develops them.

\begin{table}[htbp]\centering\small
\caption{Properties of the Fourier transform, $x(t)\leftrightarrow X(f)$, $y(t)\leftrightarrow Y(f)$.}\label{tab:appA:ftprops}
\begin{tabular}{@{}lll@{}}
\toprule
Property & Time domain & Frequency domain\\
\midrule
Linearity & $a\,x(t)+b\,y(t)$ & $a\,X(f)+b\,Y(f)$\\
Time shift & $x(t-t_0)$ & $e^{-j2\pi f t_0}X(f)$\\
Frequency shift & $x(t)\,e^{j2\pi f_0t}$ & $X(f-f_0)$\\
Modulation & $x(t)\cos 2\pi f_0 t$ & $\tfrac12X(f-f_0)+\tfrac12X(f+f_0)$\\
Scaling & $x(at)$, $a\neq0$ & $\frac{1}{|a|}X(f/a)$\\
Time reversal & $x(-t)$ & $X(-f)$\\
Conjugation & $x^*(t)$ & $X^*(-f)$\\
Real signal & $x(t)=x^*(t)$ & $X(-f)=X^*(f)$ (Hermitian symmetry)\\
Real and even & $x(t)=x(-t)$ real & $X(f)$ real and even\\
Duality & $X(t)$ & $x(-f)$\\
Convolution & $x(t)*y(t)$ & $X(f)\,Y(f)$\\
Multiplication & $x(t)\,y(t)$ & $X(f)*Y(f)$\\
Correlation & $\int x(\tau)y^*(\tau-t)\,d\tau$ & $X(f)\,Y^*(f)$\\
Differentiation & $\frac{d^n}{dt^n}x(t)$ & $(j2\pi f)^nX(f)$\\
Frequency derivative & $t\,x(t)$ & $\frac{-1}{j2\pi}\frac{d}{df}X(f)$\\
Integration & $\int_{-\infty}^{t}x(\tau)\,d\tau$ & $\frac{X(f)}{j2\pi f}+\tfrac12X(0)\delta(f)$\\
Area & $\int x(t)\,dt=X(0)$ & $\int X(f)\,df=x(0)$\\
Hilbert transform & $\hat x(t)=x(t)*\frac{1}{\pi t}$ & $-j\operatorname{sgn}(f)\,X(f)$\\
Analytic signal & $x(t)+j\hat x(t)$ & $2u(f)X(f)$\\
Parseval & $\int x(t)y^*(t)\,dt$ & $\int X(f)Y^*(f)\,df$\\
\bottomrule
\end{tabular}
\end{table}

\begin{table}[htbp]\centering\small
\caption{Fourier transform pairs. $T>0$, $a>0$, $0\le\beta\le1$.}\label{tab:appA:ftpairs}
\renewcommand{\arraystretch}{1.25}
\begin{tabular}{@{}lll@{}}
\toprule
Signal & $x(t)$ & $X(f)$\\
\midrule
Impulse & $\delta(t)$ & $1$\\
Constant & $1$ & $\delta(f)$\\
Complex exponential & $e^{j2\pi f_0t}$ & $\delta(f-f_0)$\\
Cosine & $\cos 2\pi f_0 t$ & $\tfrac12[\delta(f-f_0)+\delta(f+f_0)]$\\
Sine & $\sin 2\pi f_0 t$ & $\tfrac1{2j}[\delta(f-f_0)-\delta(f+f_0)]$\\
Rectangular pulse & $\operatorname{rect}(t/T)$ & $T\sinc(fT)$\\
Sinc pulse & $\sinc(t/T)$ & $T\operatorname{rect}(fT)$\\
Triangle & $\Lambda(t/T)$ & $T\sinc^2(fT)$\\
Gaussian & $e^{-\pi t^2/T^2}$ & $T\,e^{-\pi f^2T^2}$\\
One-sided exponential & $e^{-at}u(t)$ & $\dfrac{1}{a+j2\pi f}$\\
Two-sided exponential & $e^{-a|t|}$ & $\dfrac{2a}{a^2+(2\pi f)^2}$\\
Damped cosine & $e^{-at}\cos(2\pi f_0t)\,u(t)$ & $\dfrac{a+j2\pi f}{(a+j2\pi f)^2+(2\pi f_0)^2}$\\
Signum & $\operatorname{sgn}(t)$ & $\dfrac{1}{j\pi f}$\\
Unit step & $u(t)$ & $\tfrac12\delta(f)+\dfrac{1}{j2\pi f}$\\
Hilbert kernel & $\dfrac{1}{\pi t}$ & $-j\operatorname{sgn}(f)$\\
Impulse train (comb) & $\sum_n\delta(t-nT)$ & $\frac1T\sum_k\delta(f-k/T)$\\
Gaussian-filtered (GMSK) & $\frac{\sqrt{2\pi}B}{\sqrt{\ln2}}\,e^{-2\pi^2B^2t^2/\ln2}$ & $e^{-(\ln 2/2)(f/B)^2}$ ($B$ = 3-dB bandwidth)\\
\bottomrule
\end{tabular}
\end{table}

Two more pairs are used constantly. The \term{raised-cosine} pulse with roll-off $\beta$ and its
spectrum are
\begin{gather}
x_{\mathrm{RC}}(t)=\sinc\Big(\frac tT\Big)\frac{\cos(\pi\beta t/T)}{1-(2\beta t/T)^2}
\quad\leftrightarrow\notag\\
X_{\mathrm{RC}}(f)=\begin{cases}T,& |f|\le\frac{1-\beta}{2T},\\[2pt]
\frac T2\Big[1+\cos\Big(\frac{\pi T}{\beta}\big(|f|-\frac{1-\beta}{2T}\big)\Big)\Big],&\frac{1-\beta}{2T}<|f|\le\frac{1+\beta}{2T},\\[2pt]
0,&|f|>\frac{1+\beta}{2T},\end{cases}
\label{eq:appA:rc}
\end{gather}
and the \term{root-raised-cosine} pulse, whose spectrum is $\sqrt{T\,X_{\mathrm{RC}}(f)}$ and whose
energy is $T$, is
\begin{equation}
x_{\mathrm{RRC}}(t)=\frac{\sin\big(\pi(1-\beta)t/T\big)+4\beta(t/T)\cos\big(\pi(1+\beta)t/T\big)}{\pi(t/T)\big[1-(4\beta t/T)^2\big]}.
\label{eq:appA:rrc}
\end{equation}
Both are derived in Chapter~\ref{ch:baseband}.
The RRC pulse as written has peak value $1-\beta+4\beta/\pi$ at $t=0$; at $t=\pm T/(4\beta)$ the
formula is $0/0$ and the limit is
$\frac{\beta}{\sqrt2}\big[(1+\tfrac2\pi)\sin\tfrac{\pi}{4\beta}+(1-\tfrac2\pi)\cos\tfrac{\pi}{4\beta}\big]$.
Numerical code must handle both points; \lab{commlib.rrc\_taps} does.

\subsection{Energy, power and Parseval's theorem}

\term{Parseval's theorem} (also called Plancherel's or Rayleigh's energy theorem) says the
transform preserves inner products:
\begin{equation}
\int_{-\infty}^{\infty}x(t)\,y^*(t)\,dt=\int_{-\infty}^{\infty}X(f)\,Y^*(f)\,df,\qquad
E_x=\int|x(t)|^2dt=\int|X(f)|^2df.
\label{eq:appA:parseval}
\end{equation}
$|X(f)|^2$ is therefore the \term{energy spectral density}. Two signals with non-overlapping
spectra are orthogonal; this single observation underlies FDMA, the quadrature carriers of QAM
and the orthogonality of OFDM subcarriers. For a power signal, the power spectral density is
the limit $S_x(f)=\lim_{T_0\to\infty}\E|X_{T_0}(f)|^2/T_0$ of the energy spectrum of a
$T_0$-second segment, which for a wide-sense stationary process equals the Fourier transform of
the autocorrelation (Section~\ref{sec:appA:processes}).

For a periodic signal with period $T_0$ the Fourier series
\begin{equation}
x(t)=\sum_{k=-\infty}^{\infty}c_k\,e^{j2\pi kt/T_0},\qquad c_k=\frac1{T_0}\int_{T_0}x(t)\,e^{-j2\pi kt/T_0}\,dt
\end{equation}
gives line spectra $X(f)=\sum_k c_k\delta(f-k/T_0)$ and the power form of Parseval's theorem,
$\frac1{T_0}\int_{T_0}|x|^2dt=\sum_k|c_k|^2$. A square wave of amplitude $\pm1$ has
$|c_k|=2/(\pi |k|)$ for odd $k$, so 81\% of its power ($8/\pi^2$) is in the fundamental pair.

\subsection{Sampling identities}

Sampling $x(t)$ every $T_s=1/f_s$ seconds multiplies it by an impulse train, so by the comb pair
and the multiplication property the spectrum of the sampled signal is periodic:
\begin{equation}
x_s(t)=\sum_n x(nT_s)\,\delta(t-nT_s)\;\leftrightarrow\;X_s(f)=f_s\sum_{k}X(f-kf_s).
\label{eq:appA:sampled}
\end{equation}
The same fact, stated without impulses, is the \term{Poisson summation formula}
\begin{equation}
\sum_{n=-\infty}^{\infty}x(nT_s)\,e^{-j2\pi fnT_s}=f_s\sum_{k=-\infty}^{\infty}X(f-kf_s),
\label{eq:appA:poisson}
\end{equation}
which links the DTFT of the samples to the folded continuous spectrum. If $X(f)=0$ for
$|f|\ge f_s/2$, the copies do not overlap and the signal is recovered by ideal low-pass
interpolation,
\begin{equation}
x(t)=\sum_{n}x(nT_s)\,\sinc\!\left(\frac{t-nT_s}{T_s}\right)
\label{eq:appA:interp}
\end{equation}
(the Whittaker--Kotelnikov--Shannon theorem, Chapter~\ref{ch:pcm}). Applied to a pulse rather
than a signal, \eqref{eq:appA:poisson} is the Nyquist zero-ISI criterion of
Chapter~\ref{ch:baseband}: $p(nT)=\delta[n]$ if and only if $\sum_kP(f-k/T)=T$.

\paragraph{Bandpass sampling.} A real signal occupying $f_L\le|f|\le f_H$ (bandwidth
$B=f_H-f_L$) can be sampled without aliasing at any rate satisfying
\begin{equation}
\frac{2f_H}{n}\le f_s\le\frac{2f_L}{n-1},\qquad n=1,2,\dots,\left\lfloor f_H/B\right\rfloor,
\label{eq:appA:bandpass}
\end{equation}
the largest $n$ giving the lowest usable rate, which approaches $2B$ when $f_H$ is an integer
multiple of $B$. Complex (IQ) sampling of the complex envelope needs only $f_s\ge B$
(Chapter~\ref{ch:sdr}).

\subsection{The discrete-time Fourier transform}

For a sequence $x[n]$ the DTFT and its inverse are
\begin{equation}
X(e^{j\omega})=\sum_{n=-\infty}^{\infty}x[n]\,e^{-j\omega n},\qquad
x[n]=\frac1{2\pi}\int_{-\pi}^{\pi}X(e^{j\omega})\,e^{j\omega n}\,d\omega,
\label{eq:appA:dtft}
\end{equation}
periodic in $\omega$ with period $2\pi$. When $x[n]=x_c(nT_s)$, the frequency variables are
related by $\omega=2\pi f/f_s$ and $X(e^{j\omega})=f_s\sum_kX_c(f-kf_s)$ by~\eqref{eq:appA:poisson}.
The engineer's normalised frequency $\nu=f/f_s$ in cycles per sample (used throughout
\texttt{commlib}) is $\omega/2\pi$.

\begin{table}[htbp]\centering\small
\caption{DTFT properties and pairs. $|a|<1$ where stated; $\omega$ is in radians per sample.}\label{tab:appA:dtft}
\begin{tabular}{@{}lll@{}}
\toprule
 & $x[n]$ & $X(e^{j\omega})$\\
\midrule
Shift & $x[n-n_0]$ & $e^{-j\omega n_0}X(e^{j\omega})$\\
Modulation & $e^{j\omega_0n}x[n]$ & $X(e^{j(\omega-\omega_0)})$\\
Convolution & $x[n]*y[n]$ & $X(e^{j\omega})Y(e^{j\omega})$\\
Multiplication & $x[n]\,y[n]$ & $\frac1{2\pi}\int_{-\pi}^{\pi}X(e^{j\theta})Y(e^{j(\omega-\theta)})\,d\theta$\\
Parseval & $\sum_n|x[n]|^2$ & $\frac1{2\pi}\int_{-\pi}^{\pi}|X(e^{j\omega})|^2d\omega$\\
Upsampling by $L$ & $x[n/L]$ ($n$ multiple of $L$, else 0) & $X(e^{j\omega L})$\\
Downsampling by $M$ & $x[nM]$ & $\frac1M\sum_{k=0}^{M-1}X(e^{j(\omega-2\pi k)/M})$\\
\midrule
Impulse & $\delta[n]$ & $1$\\
Exponential & $a^n u[n]$ & $\dfrac{1}{1-ae^{-j\omega}}$\\
Rectangular window & $1$ for $0\le n\le N-1$ & $e^{-j\omega(N-1)/2}\,\dfrac{\sin(N\omega/2)}{\sin(\omega/2)}$\\
Ideal low-pass & $\dfrac{\omega_c}{\pi}\sinc\!\left(\dfrac{\omega_c n}{\pi}\right)$ & $1$ for $|\omega|<\omega_c$, else 0 (per period)\\
Complex exponential & $e^{j\omega_0 n}$ & $2\pi\sum_k\delta(\omega-\omega_0-2\pi k)$\\
\bottomrule
\end{tabular}
\end{table}

The ratio $\sin(N\omega/2)/(N\sin(\omega/2))$ in Table~\ref{tab:appA:dtft} is the
\term{Dirichlet kernel} or periodic sinc. It is the spectrum of every finite observation, the
array factor of every uniform linear array (Chapter~\ref{ch:mimo}) and the response of every
moving-average and CIC filter (Chapter~\ref{ch:dsp}).

\subsection{The discrete Fourier transform}

The $N$-point DFT and its inverse are
\begin{equation}
X[k]=\sum_{n=0}^{N-1}x[n]\,e^{-j2\pi kn/N},\qquad
x[n]=\frac1N\sum_{k=0}^{N-1}X[k]\,e^{j2\pi kn/N}.
\label{eq:appA:dft}
\end{equation}
$X[k]$ samples the DTFT of the $N$-point block at $\omega_k=2\pi k/N$. The DFT treats its input as
one period of a periodic sequence, which explains all its properties
(Table~\ref{tab:appA:dft}): shifts are circular, products of DFTs correspond to
\emph{circular} convolution, and a block that does not contain an integer number of cycles
leaks into every bin. The FFT computes~\eqref{eq:appA:dft} in $\tfrac N2\log_2N$ complex
multiplications for radix-2 $N$, instead of $N^2$.

\begin{table}[htbp]\centering\small
\caption{DFT properties ($N$-point; indices modulo $N$; $W_N=e^{-j2\pi/N}$).}\label{tab:appA:dft}
\begin{tabular}{@{}lll@{}}
\toprule
Property & $x[n]$ & $X[k]$\\
\midrule
Circular shift & $x[(n-m)_N]$ & $W_N^{km}X[k]$\\
Modulation & $W_N^{-ln}x[n]$ & $X[(k-l)_N]$\\
Circular convolution & $x[n]\circledast y[n]=\sum_m x[m]\,y[(n-m)_N]$ & $X[k]\,Y[k]$\\
Multiplication & $x[n]\,y[n]$ & $\frac1N X[k]\circledast Y[k]$\\
Conjugation & $x^*[n]$ & $X^*[(-k)_N]$\\
Real input & $x[n]$ real & $X[(-k)_N]=X^*[k]$\\
Parseval & $\sum_n|x[n]|^2$ & $\frac1N\sum_k|X[k]|^2$\\
Duality & $X[n]$ & $N\,x[(-k)_N]$\\
\midrule
Impulse & $\delta[n]$ & $1$\\
Constant & $1$ & $N\delta[k]$\\
On-bin tone & $e^{j2\pi k_0n/N}$ & $N\delta[k-k_0]$\\
Off-bin tone & $e^{j2\pi(k_0+\epsilon)n/N}$ & $e^{j\pi(N-1)(k_0+\epsilon-k)/N}\dfrac{\sin\pi(k_0+\epsilon-k)}{\sin(\pi(k_0+\epsilon-k)/N)}$\\
\bottomrule
\end{tabular}
\end{table}

\begin{keyidea}[Why a cyclic prefix works]
Linear convolution with a channel of length $L$ becomes circular convolution if each block is
preceded by a copy of its last $L-1$ samples. By the circular-convolution row of
Table~\ref{tab:appA:dft}, the receiver's DFT then sees $Y[k]=H[k]X[k]$: one complex
multiplication per subcarrier. That is the whole trick of OFDM (Chapter~\ref{ch:ofdm}), and it
is also the frequency-domain equaliser of single-carrier systems and the overlap-save fast
convolution of Chapter~\ref{ch:dsp}. Section~\ref{sec:appA:circulant} states the same fact in
matrix language.
\end{keyidea}

\paragraph{Resolution and zero padding.} An $N$-point record at rate $f_s$ gives bins
$\Delta f=f_s/N$ apart; the ability to separate two tones is set by the record length
$T_0=N/f_s$ and the window, not by zero padding, which only interpolates the same DTFT more
finely. Table~\ref{tab:appA:windows} gives the classic window figures of merit after
Harris (1978), in bins of $f_s/N$.

\begin{table}[htbp]\centering\small
\caption{Common windows. Main-lobe width is null to null; ENBW is the equivalent noise
bandwidth; scalloping loss is the worst-case gain loss for a tone midway between bins.}\label{tab:appA:windows}
\begin{tabular}{@{}lccccc@{}}
\toprule
Window & Main lobe & Peak sidelobe & Sidelobe decay & ENBW & Scalloping\\
 & (bins) & (dB) & (dB/octave) & (bins) & loss (dB)\\
\midrule
Rectangular & 2 & $-13.3$ & $-6$ & 1.00 & 3.92\\
Hann & 4 & $-31.5$ & $-18$ & 1.50 & 1.42\\
Hamming & 4 & $-42.7$ & $-6$ & 1.36 & 1.78\\
Blackman & 6 & $-58.1$ & $-18$ & 1.73 & 1.10\\
Blackman--Harris (4-term) & 8 & $-92$ & $-6$ & 2.00 & 0.83\\
\bottomrule
\end{tabular}
\end{table}

\subsection{The $z$-transform}

The two-sided $z$-transform
\begin{equation}
X(z)=\sum_{n=-\infty}^{\infty}x[n]\,z^{-n}
\label{eq:appA:z}
\end{equation}
converges in an annulus called the region of convergence (ROC); on the unit circle it reduces
to the DTFT, $X(e^{j\omega})=X(z)|_{z=e^{j\omega}}$, provided the ROC contains $|z|=1$. A causal
LTI system is stable if and only if all poles of its transfer function lie strictly inside the
unit circle. Table~\ref{tab:appA:z} lists the pairs used in Chapters~\ref{ch:dsp},
\ref{ch:sync} and~\ref{ch:equalization}.

\begin{table}[htbp]\centering\small
\caption{$z$-transform pairs and properties.}\label{tab:appA:z}
\begin{tabular}{@{}lll@{}}
\toprule
$x[n]$ & $X(z)$ & ROC\\
\midrule
$\delta[n-m]$ & $z^{-m}$ & all $z$ (except $0$ or $\infty$)\\
$a^n u[n]$ & $\dfrac{1}{1-az^{-1}}$ & $|z|>|a|$\\
$-a^n u[-n-1]$ & $\dfrac{1}{1-az^{-1}}$ & $|z|<|a|$\\
$n\,a^n u[n]$ & $\dfrac{az^{-1}}{(1-az^{-1})^2}$ & $|z|>|a|$\\
$r^n\cos(\omega_0n)\,u[n]$ & $\dfrac{1-r\cos\omega_0\,z^{-1}}{1-2r\cos\omega_0\,z^{-1}+r^2z^{-2}}$ & $|z|>r$\\
$r^n\sin(\omega_0n)\,u[n]$ & $\dfrac{r\sin\omega_0\,z^{-1}}{1-2r\cos\omega_0\,z^{-1}+r^2z^{-2}}$ & $|z|>r$\\
$u[n]-u[n-N]$ & $\dfrac{1-z^{-N}}{1-z^{-1}}$ & $z\ne0$\\
\midrule
$x[n-m]$ & $z^{-m}X(z)$ & same\\
$a^nx[n]$ & $X(z/a)$ & scaled by $|a|$\\
$x[-n]$ & $X(1/z)$ & inverted\\
$x^*[n]$ & $X^*(z^*)$ & same\\
$n\,x[n]$ & $-z\,\dfrac{dX}{dz}$ & same\\
$x[n]*y[n]$ & $X(z)Y(z)$ & at least the intersection\\
$\sum_{m\le n}x[m]$ (accumulator) & $\dfrac{X(z)}{1-z^{-1}}$ & includes $|z|>1$\\
Initial value (causal $x$) & $x[0]=\lim_{z\to\infty}X(z)$ & \\
Final value & $\lim_{n\to\infty}x[n]=\lim_{z\to1}(1-z^{-1})X(z)$ & if the limit exists\\
\bottomrule
\end{tabular}
\end{table}

The accumulator $1/(1-z^{-1})$ and the final-value theorem together explain why a type-2
phase-locked loop (one with an integrator in its loop filter) tracks a frequency step with zero
steady-state phase error (Chapter~\ref{ch:sync}). Analog prototypes are carried into discrete
time by the \term{bilinear transform}
\begin{equation}
s=\frac{2}{T_s}\,\frac{1-z^{-1}}{1+z^{-1}},\qquad \Omega=\frac{2}{T_s}\tan\frac{\omega}{2},
\label{eq:appA:bilinear}
\end{equation}
which maps the left half-plane onto the unit disc and warps frequency; designs are pre-warped
with the second relation (Chapter~\ref{ch:dsp}). The Laplace pairs $e^{-at}u(t)\leftrightarrow
1/(s+a)$ and $\delta(t)\leftrightarrow1$ are the continuous counterparts of the first two rows.

%======================================================================
\section{Special functions}
\label{sec:appA:special}
%======================================================================

\subsection{Pulses and kernels}

The \term{rectangle} $\operatorname{rect}(t)$ equals 1 for $|t|<\tfrac12$ and 0 outside; the
\term{triangle} $\Lambda(t)=(1-|t|)^+=\operatorname{rect}*\operatorname{rect}$; the sinc
$\sinc x=\sin(\pi x)/(\pi x)$ has unit area, unit energy ($\int\sinc^2=1$), zeros at the
non-zero integers and sidelobes that decay as $1/(\pi|x|)$: the first is at $x\approx1.43$
with height $-0.217$ ($-13.3$\,dB). Its square, the spectrum of a rectangular NRZ pulse, has its
first null at $f=1/T$ and contains 90.3\% of its power within $|f|<1/T$. The Dirichlet kernel
$D_N(\omega)=\sin(N\omega/2)/(N\sin(\omega/2))$ is the periodic counterpart and approaches
$\sinc(N\omega/2\pi)$ for small $\omega$.

\subsection{The Gaussian Q-function and the error function}

The \term{Q-function} is the tail of the standard normal distribution,
\begin{equation}
\Q(x)=\frac{1}{\sqrt{2\pi}}\int_x^{\infty}e^{-u^2/2}\,du=\tfrac12\operatorname{erfc}\!\left(\frac{x}{\sqrt2}\right),
\label{eq:appA:q}
\end{equation}
where $\operatorname{erfc}(x)=\frac{2}{\sqrt\pi}\int_x^\infty e^{-u^2}du=2\Q(\sqrt2\,x)$ is the
complementary error function,
and $\Q(-x)=1-\Q(x)$, $\Q(0)=\tfrac12$ and $\frac{d}{dx}\Q(x)=-\phi(x)$, where
$\phi(x)=e^{-x^2/2}/\sqrt{2\pi}$ is the standard normal density. Almost every error rate in
Chapters~\ref{ch:baseband}--\ref{ch:modulation} is a Q-function of a distance divided by a noise
standard deviation: binary antipodal signalling gives $P_b=\Q(\sqrt{2E_b/N_0})$. Numerical
libraries provide \texttt{erfc} (and \texttt{scipy.special.ndtr}, \texttt{scipy.stats.norm.sf});
\texttt{commlib.qfunc} wraps it.

\paragraph{Craig's form.} For $x\ge0$, Craig (1991) showed that
\begin{equation}
\Q(x)=\frac1\pi\int_0^{\pi/2}\exp\!\left(-\frac{x^2}{2\sin^2\theta}\right)d\theta,\qquad
\Q^2(x)=\frac1\pi\int_0^{\pi/4}\exp\!\left(-\frac{x^2}{2\sin^2\theta}\right)d\theta.
\label{eq:appA:craig}
\end{equation}
The argument now appears only in the exponent and the limits are finite, so averaging over a
random SNR becomes an integral of the SNR's moment-generating function
(Section~\ref{sec:appA:mgfmethod}). This is the key to almost every closed-form fading error
rate in Chapters~\ref{ch:channels} and~\ref{ch:mimo}.

\paragraph{Bounds and approximations.} For $x>0$,
\begin{equation}
\frac{x}{1+x^2}\,\phi(x)<\Q(x)<\frac{\phi(x)}{x},\qquad
\Q(x)\le\tfrac12e^{-x^2/2},\qquad
\Q(x)\approx\tfrac1{12}e^{-x^2/2}+\tfrac14e^{-2x^2/3}.
\label{eq:appA:qbounds}
\end{equation}
The first pair is tight for large $x$ (the ratio of the bounds tends to 1); the second is the
Chernoff bound, loose by a factor that grows like $x$ but convenient because it has the exact
exponent and multiplies nicely; the third, from Chiani, Dardari and Simon (2003), is a sum of
exponentials that is accurate to about 1\,dB over the whole range and can be averaged over fading
in closed form. Figure~\ref{fig:appA_qfunc} compares them, and Table~\ref{tab:appA:qvalues}
gives the values engineers memorise.

\mdcfig{appA_qfunc}{The Gaussian Q-function and its standard bounds. Left: values on a log
scale. Right: the error of each approximation relative to the exact value, in dB. The Chernoff
bound has the right exponent but is about 2--5\,dB pessimistic in the useful range; the
asymptotic bounds $\phi(x)/x$ and $\phi(x)(1-x^{-2})/x$ bracket $\Q(x)$ tightly once $x>3$
(error rates below $10^{-3}$).}

\begin{table}[htbp]\centering\small
\caption{The Q-function and its inverse. The last column is the $E_b/N_0$ at which binary
antipodal signalling ($P_b=\Q(\sqrt{2E_b/N_0})$) reaches the error rate in the first column.}\label{tab:appA:qvalues}
\begin{tabular}{@{}cc@{\qquad}ccc@{}}
\toprule
$x$ & $\Q(x)$ & $\Q(x)$ & $x=\Q^{-1}(\cdot)$ & BPSK $E_b/N_0$ (dB)\\
\midrule
0.0 & $5.000\times10^{-1}$ & $10^{-1}$ & 1.2816 & $-0.86$\\
1.0 & $1.587\times10^{-1}$ & $10^{-2}$ & 2.3263 & 4.32\\
2.0 & $2.275\times10^{-2}$ & $10^{-3}$ & 3.0902 & 6.79\\
3.0 & $1.350\times10^{-3}$ & $10^{-4}$ & 3.7190 & 8.40\\
3.5 & $2.326\times10^{-4}$ & $10^{-5}$ & 4.2649 & 9.59\\
4.0 & $3.167\times10^{-5}$ & $10^{-6}$ & 4.7534 & 10.53\\
4.5 & $3.398\times10^{-6}$ & $10^{-7}$ & 5.1993 & 11.31\\
5.0 & $2.867\times10^{-7}$ & $10^{-8}$ & 5.6120 & 11.97\\
5.5 & $1.899\times10^{-8}$ & $10^{-9}$ & 5.9978 & 12.55\\
6.0 & $9.866\times10^{-10}$ & $10^{-10}$ & 6.3613 & 13.06\\
6.5 & $4.016\times10^{-11}$ & $10^{-12}$ & 7.0345 & 13.93\\
7.0 & $1.280\times10^{-12}$ & $10^{-15}$ & 7.9413 & 14.99\\
\bottomrule
\end{tabular}
\end{table}

\begin{inpractice}[The ``Q-factor'' of optical and SerDes engineers]
Optical-link and SerDes engineers quote a link's quality as a \emph{Q-factor}, the eye opening
divided by the sum of the noise standard deviations on the two rails, so that $\text{BER}=\Q(Q)$
for binary signalling. Read Table~\ref{tab:appA:qvalues} sideways: $Q=7.03$ (often written
$20\log_{10}Q=16.9$\,dB) is the classic $10^{-12}$ target, and $Q\approx3.1$ ($10^{-3}$) is the
pre-FEC threshold of many modern links that rely on Reed--Solomon or LDPC codes to clean up the
rest (Chapters~\ref{ch:baseband} and~\ref{ch:wireline}).
\end{inpractice}

\subsection{Marcum's Q-function}

The generalised \term{Marcum Q-function} of order $M$ is
\begin{equation}
\Q_M(a,b)=\int_b^\infty x\left(\frac xa\right)^{M-1}\exp\!\left(-\frac{x^2+a^2}{2}\right)I_{M-1}(ax)\,dx,
\label{eq:appA:marcum}
\end{equation}
where $I_\nu$ is the modified Bessel function below. $\Q_1(a,b)$ is the probability that a
Rician envelope with line-of-sight amplitude $a$ and unit-variance Gaussian components exceeds
$b$; equivalently $\Q_M(a,b)=\Pr\{\chi'^2_{2M}(a^2)>b^2\}$, the tail of a noncentral chi-square
variable with $2M$ degrees of freedom. Special values: $\Q_1(a,0)=1$, $\Q_1(0,b)=e^{-b^2/2}$
(Rayleigh), and $\Q_1(a,b)\approx\Q(b-a)$ when $a$ and $b$ are both large. It gives the
detection probability of a noncoherent radar or energy detector, the error rate of noncoherent
FSK and DPSK with diversity, and the outage probability of Rician links
(Chapters~\ref{ch:modulation} and~\ref{ch:channels}). In Python,
$\Q_M(a,b)=$ \texttt{scipy.stats.ncx2.sf(b**2, 2*M, a**2)}.

\subsection{Bessel functions}

The Bessel function of the first kind of integer order $n$ and the modified Bessel function
of the first kind are
\begin{equation}
J_n(x)=\frac1\pi\int_0^\pi\cos(n\theta-x\sin\theta)\,d\theta,\qquad
I_n(x)=\frac1\pi\int_0^\pi e^{x\cos\theta}\cos(n\theta)\,d\theta.
\label{eq:appA:bessel}
\end{equation}

\paragraph{$J_n$ and FM.} The Jacobi--Anger expansion
\begin{equation}
e^{j\beta\sin\theta}=\sum_{n=-\infty}^{\infty}J_n(\beta)\,e^{jn\theta}
\label{eq:appA:jacobi}
\end{equation}
says that a carrier phase-modulated by a sinusoid with index $\beta$ has spectral lines at
$f_c+nf_m$ with amplitudes $J_n(\beta)$ (Chapter~\ref{ch:analog}). Useful facts:
$J_{-n}(x)=(-1)^nJ_n(x)$; $\sum_nJ_n^2(x)=1$ (power is conserved); for small $x$,
$J_0(x)\approx1-x^2/4$ and $J_1(x)\approx x/2$; and the carrier vanishes at the zeros of $J_0$,
$\beta=2.405$, 5.520, 8.654, 11.79, which is how FM deviation meters are calibrated (the
``Bessel null'' method). Carson's rule follows from the fact that $J_n(\beta)$ is small for
$|n|>\beta+1$.

\paragraph{$I_0$ and noncoherent detection.} $I_0(x)=\frac1{2\pi}\int_0^{2\pi}e^{x\cos\theta}d\theta$
is what remains when a Gaussian likelihood is averaged over a uniformly distributed carrier
phase, so it appears in the Rician density, the noncoherent receiver (Chapter~\ref{ch:modulation}),
the von Mises phase distribution of a PLL (Chapter~\ref{ch:sync}) and the Kaiser window
(Chapter~\ref{ch:dsp}). Series and asymptote:
\begin{equation}
I_0(x)=\sum_{k=0}^{\infty}\frac{(x/2)^{2k}}{(k!)^2}\approx1+\frac{x^2}{4}\;(x\ll1),\qquad
I_0(x)\approx\frac{e^x}{\sqrt{2\pi x}}\;(x\gg1).
\label{eq:appA:i0}
\end{equation}
Large arguments overflow in floating point; use the exponentially scaled form
$I_0(x)e^{-x}$ (\texttt{scipy.special.i0e}) and work with logarithms. The modified Bessel function
of the second kind $K_\nu$ appears in the density of a product of two Gaussians and in the
$K$-distribution used for composite fading and sea clutter.

\mdcfig{appA_bessel_marcum}{Left: Bessel functions $J_0$--$J_3$, the amplitudes of FM sidebands
versus modulation index. Centre: modified Bessel functions $I_0$ and $I_1$ with the large-argument
asymptote. Right: Marcum $\Q_1(a,b)$, the probability that a Rician envelope with line-of-sight
amplitude $a$ (unit-variance components) exceeds a threshold $b$; $a=0$ is the Rayleigh case
$e^{-b^2/2}$.}

\subsection{Gamma, incomplete gamma and related functions}

The \term{gamma function} $\Gamma(z)=\int_0^\infty t^{z-1}e^{-t}dt$ extends the factorial:
$\Gamma(n)=(n-1)!$, $\Gamma(z+1)=z\Gamma(z)$, $\Gamma(\tfrac12)=\sqrt\pi$. Stirling's
approximation is $\ln n!\approx n\ln n-n+\tfrac12\ln(2\pi n)$. The lower and upper incomplete
gamma functions split the defining integral at $x$,
\begin{equation}
\gamma(s,x)=\int_0^x t^{s-1}e^{-t}dt,\qquad \Gamma(s,x)=\int_x^\infty t^{s-1}e^{-t}dt,\qquad
P(s,x)=\frac{\gamma(s,x)}{\Gamma(s)},
\label{eq:appA:incgamma}
\end{equation}
and the regularised form $P(s,x)$ (\texttt{scipy.special.gammainc}) is the CDF of a gamma variable.
It therefore gives, without further work, the CDF of chi-square ($P(k/2,x/2)$), of the power of a
Nakagami-$m$ signal, of the output SNR of $L$-branch maximal-ratio combining in Rayleigh fading
($P(L,x/\bar\gamma)$) and of the Poisson distribution ($\Pr\{N\le k\}=1-P(k+1,\lambda)$). The
beta function is $B(a,b)=\Gamma(a)\Gamma(b)/\Gamma(a+b)$.

The \term{exponential integral} $E_1(x)=\int_x^\infty e^{-t}t^{-1}dt$ gives the ergodic capacity
of a Rayleigh channel with receiver CSI,
\begin{equation}
\E\big[\log_2(1+\bar\gamma|h|^2)\big]=\log_2(e)\;e^{1/\bar\gamma}\,E_1(1/\bar\gamma),
\label{eq:appA:ergodic}
\end{equation}
which at $\bar\gamma=10$\,dB is 2.91\,b/s/Hz against 3.46 for an unfaded AWGN channel at the same
mean SNR (Chapters~\ref{ch:infotheory} and~\ref{ch:mimo}).

%======================================================================
\section{Probability and random variables}
\label{sec:appA:probability}
%======================================================================

\subsection{Essentials}

A random variable $X$ has CDF $F_X(x)=\Pr\{X\le x\}$ and density $f_X(x)=F_X'(x)$. Its mean is
$\E X=\int xf_X(x)dx$, its variance $\operatorname{Var}X=\E X^2-(\E X)^2$, and its
\term{moment-generating function} (MGF) and characteristic function are
\begin{equation}
M_X(s)=\E\big[e^{sX}\big],\qquad \Phi_X(\omega)=\E\big[e^{j\omega X}\big]=M_X(j\omega),
\end{equation}
with $\E X^n=M_X^{(n)}(0)$. The MGF of a sum of independent variables is the product of their
MGFs, which is why the MGF is the natural tool for diversity combining. Other workhorses:

\begin{itemize}[leftmargin=1.5em,itemsep=2pt]
\item \emph{Bayes' rule}: $\Pr\{A|B\}=\Pr\{B|A\}\Pr\{A\}/\Pr\{B\}$; with densities,
$p(x|y)=p(y|x)p(x)/p(y)$. MAP detection maximises $p(y|x)p(x)$; ML detection maximises $p(y|x)$.
\item \emph{Total probability}: $\Pr\{B\}=\sum_i\Pr\{B|A_i\}\Pr\{A_i\}$ over a partition $\{A_i\}$.
\item \emph{Change of variables}: if $Y=g(X)$ with roots $x_i$ of $g(x_i)=y$, then
$f_Y(y)=\sum_if_X(x_i)/|g'(x_i)|$. For vectors, divide by $|\det\vect{J}|$, the Jacobian.
\item \emph{Sums}: the density of $X+Y$ for independent $X,Y$ is $f_X*f_Y$.
\item \emph{Inequalities}: Markov, $\Pr\{X\ge a\}\le\E X/a$ for $X\ge0$; Chebyshev,
$\Pr\{|X-\mu|\ge k\sigma\}\le1/k^2$; Chernoff, $\Pr\{X\ge a\}\le\min_{s>0}e^{-sa}M_X(s)$;
Jensen, $g(\E X)\le\E g(X)$ for convex $g$.
\item \emph{Limit theorems}: the sample mean converges to $\E X$ (law of large numbers) and the
normalised sum of many independent finite-variance terms tends to a Gaussian (central limit
theorem, Chapter~\ref{ch:noise}).
\item \emph{Monte Carlo accuracy}: estimating a probability $p$ from $N$ trials gives relative
standard deviation $\sqrt{(1-p)/(Np)}\approx1/\sqrt{\text{number of errors}}$; 100 errors give
about $\pm20\%$ at 95\% confidence. This is the stopping rule of \texttt{labkit.ber\_mc}.
\end{itemize}

\subsection{A table of distributions}

Tables~\ref{tab:appA:dists} and~\ref{tab:appA:moments} summarise the continuous distributions of
communications, and Table~\ref{tab:appA:discrete} the discrete ones. The Rayleigh, Rice and
Nakagami entries describe the \emph{envelope} of a fading channel normalised to mean power
$\Omega=\E R^2$; the power $R^2$ is then exponential, noncentral chi-square or gamma
respectively. Figures~\ref{fig:appA_fading_pdfs} and~\ref{fig:appA_chi2} plot the shapes.

\begin{table}[htbp]\centering\footnotesize
\caption{Continuous distributions: density and CDF. $\Q$ is the Gaussian tail function,
$\Q_M$ Marcum's function, $P(s,x)$ the regularised lower incomplete gamma function.}\label{tab:appA:dists}
\renewcommand{\arraystretch}{1.55}
\begin{tabular}{@{}lll@{}}
\toprule
Distribution & Density $f(x)$ & CDF $F(x)$\\
\midrule
Gaussian $\mathcal N(\mu,\sigma^2)$ & $\frac{1}{\sqrt{2\pi\sigma^2}}e^{-(x-\mu)^2/2\sigma^2}$ & $1-\Q\big(\frac{x-\mu}{\sigma}\big)$\\
Complex Gaussian $\mathcal{CN}(0,\sigma^2)$ & $\frac{1}{\pi\sigma^2}e^{-|z|^2/\sigma^2}$, $z\in\C$ & ---\\
Uniform $\mathcal U(a,b)$ & $\frac1{b-a}$, $a\le x\le b$ & $\frac{x-a}{b-a}$\\
Exponential (rate $\lambda$) & $\lambda e^{-\lambda x}$, $x\ge0$ & $1-e^{-\lambda x}$\\
Rayleigh ($\sigma$) & $\frac{x}{\sigma^2}e^{-x^2/2\sigma^2}$, $x\ge0$ & $1-e^{-x^2/2\sigma^2}$\\
Rice ($\nu,\sigma$) & $\frac{x}{\sigma^2}e^{-(x^2+\nu^2)/2\sigma^2}I_0\big(\frac{x\nu}{\sigma^2}\big)$ & $1-\Q_1\big(\frac\nu\sigma,\frac x\sigma\big)$\\
Nakagami-$m$ ($m\ge\frac12$, $\Omega$) & $\frac{2m^m x^{2m-1}}{\Gamma(m)\Omega^m}e^{-mx^2/\Omega}$ & $P\big(m,\frac{mx^2}{\Omega}\big)$\\
Gamma ($k,\theta$) & $\frac{x^{k-1}e^{-x/\theta}}{\Gamma(k)\theta^k}$, $x\ge0$ & $P(k,x/\theta)$\\
Chi-square $\chi^2_k$ & $\frac{x^{k/2-1}e^{-x/2}}{2^{k/2}\Gamma(k/2)}$, $x\ge0$ & $P\big(\frac k2,\frac x2\big)$\\
Noncentral $\chi'^2_k(\lambda)$ & $\frac12e^{-(x+\lambda)/2}\big(\frac x\lambda\big)^{k/4-1/2}I_{k/2-1}(\sqrt{\lambda x})$ & $1-\Q_{k/2}(\sqrt\lambda,\sqrt x)$\\
Log-normal ($\mu,\sigma$) & $\frac{1}{x\sigma\sqrt{2\pi}}e^{-(\ln x-\mu)^2/2\sigma^2}$, $x>0$ & $1-\Q\big(\frac{\ln x-\mu}{\sigma}\big)$\\
\bottomrule
\end{tabular}
\end{table}

\begin{table}[htbp]\centering\footnotesize
\caption{Continuous distributions: moments and MGF $M(s)=\E e^{sX}$. $L_{1/2}$ is a Laguerre
polynomial; ``---'' means no simple closed form.}\label{tab:appA:moments}
\renewcommand{\arraystretch}{1.55}
\begin{tabular}{@{}llll@{}}
\toprule
Distribution & Mean & Variance & MGF $M(s)$\\
\midrule
Gaussian & $\mu$ & $\sigma^2$ & $e^{\mu s+\sigma^2s^2/2}$\\
Complex Gaussian & $0$ & $\E|z|^2=\sigma^2$ & $|z|^2$ is exponential, mean $\sigma^2$\\
Uniform & $\frac{a+b}2$ & $\frac{(b-a)^2}{12}$ & $\frac{e^{sb}-e^{sa}}{s(b-a)}$\\
Exponential & $\frac1\lambda$ & $\frac1{\lambda^2}$ & $\frac{\lambda}{\lambda-s}$, $s<\lambda$\\
Rayleigh & $\sigma\sqrt{\pi/2}$ & $\big(2-\frac\pi2\big)\sigma^2$ & $\E R^2=2\sigma^2$\\
Rice & $\sigma\sqrt{\frac\pi2}\,L_{1/2}\big(-\frac{\nu^2}{2\sigma^2}\big)$ & $\nu^2+2\sigma^2-(\E R)^2$ & $\E R^2=\nu^2+2\sigma^2$, $K=\frac{\nu^2}{2\sigma^2}$\\
Nakagami-$m$ & $\frac{\Gamma(m+\frac12)}{\Gamma(m)}\sqrt{\frac\Omega m}$ & $\Omega-(\E R)^2$ & $\E R^2=\Omega$; $R^2\sim$ Gamma$(m,\Omega/m)$\\
Gamma & $k\theta$ & $k\theta^2$ & $(1-\theta s)^{-k}$\\
Chi-square & $k$ & $2k$ & $(1-2s)^{-k/2}$\\
Noncentral $\chi'^2$ & $k+\lambda$ & $2(k+2\lambda)$ & $\frac{\exp(\lambda s/(1-2s))}{(1-2s)^{k/2}}$\\
Log-normal & $e^{\mu+\sigma^2/2}$ & $(e^{\sigma^2}-1)e^{2\mu+\sigma^2}$ & ---\\
\bottomrule
\end{tabular}
\end{table}

\begin{table}[htbp]\centering\small
\caption{Discrete distributions ($0<p<1$, $\lambda>0$).}\label{tab:appA:discrete}
\renewcommand{\arraystretch}{1.35}
\begin{tabular}{@{}lllll@{}}
\toprule
Distribution & $\Pr\{X=k\}$ & Mean & Variance & MGF\\
\midrule
Bernoulli & $p^k(1-p)^{1-k}$, $k\in\{0,1\}$ & $p$ & $p(1-p)$ & $1-p+pe^s$\\
Binomial $(n,p)$ & $\binom nk p^k(1-p)^{n-k}$ & $np$ & $np(1-p)$ & $(1-p+pe^s)^n$\\
Geometric (trials to success) & $p(1-p)^{k-1}$, $k\ge1$ & $\frac1p$ & $\frac{1-p}{p^2}$ & $\frac{pe^s}{1-(1-p)e^s}$\\
Poisson $(\lambda)$ & $\frac{\lambda^ke^{-\lambda}}{k!}$ & $\lambda$ & $\lambda$ & $e^{\lambda(e^s-1)}$\\
\bottomrule
\end{tabular}
\end{table}

\paragraph{How the distributions connect.}
\begin{itemize}[leftmargin=1.5em,itemsep=2pt]
\item If $Z_1,\dots,Z_k$ are independent $\mathcal N(0,1)$, $\sum Z_i^2\sim\chi^2_k$; with means
$\mu_i$, $\sum Z_i^2\sim\chi'^2_k(\lambda)$ with $\lambda=\sum\mu_i^2$.
\item If $z\sim\mathcal{CN}(0,\sigma^2)$, then $|z|$ is Rayleigh with parameter $\sigma/\sqrt2$,
$|z|^2$ is exponential with mean $\sigma^2$ ($=\tfrac{\sigma^2}{2}\chi^2_2$), and $\arg z$ is
uniform. Adding a constant (line of sight) makes $|z|$ Rician.
\item The sum of $L$ independent exponentials with mean $\bar\gamma$ is Gamma$(L,\bar\gamma)$
(Erlang): the SNR after $L$-branch maximal-ratio combining in i.i.d.\ Rayleigh fading.
\item Nakagami-$m$ with $m=1$ is Rayleigh and with $m=\tfrac12$ is one-sided Gaussian; a Rician
channel with factor $K$ is approximated by Nakagami with $m=(K+1)^2/(2K+1)$.
\item Shadowing quoted as ``$\sigma=8$\,dB'' is log-normal with $\sigma_{\ln}=\frac{\ln10}{10}\sigma_{\mathrm{dB}}=0.230\,\sigma_{\mathrm{dB}}$.
\item Poisson arrivals at rate $\lambda$ have exponential inter-arrival times with mean $1/\lambda$;
the number in a window of length $t$ is Poisson$(\lambda t)$ (Chapter~\ref{ch:multipleaccess}).
\end{itemize}

\mdcfig{appA_fading_pdfs}{Envelope densities of the standard fading models at unit mean power.
Left: Rician fading from Rayleigh ($K=0$) to a strong line of sight ($K=10$). Right: Nakagami-$m$,
from worse than Rayleigh ($m=\tfrac12$) to nearly constant ($m=10$). The behaviour near $r=0$,
which sets the diversity order, goes as $r^{2m-1}$.}

\mdcfig{appA_chi2}{Left: central chi-square densities with $k$ degrees of freedom (the energy
of $k$ unit-variance Gaussian components); $k=2$ is the exponential power of Rayleigh fading.
Right: noncentral chi-square with two degrees of freedom, the output of an energy detector or the
power of a Rician signal, for increasing signal energy $\lambda$.}

\subsection{Error probability in fading: the MGF method}
\label{sec:appA:mgfmethod}

Suppose the instantaneous SNR $\gamma$ is random and the conditional error rate is
$\Q(\sqrt{2g\gamma})$, with $g=1$ for BPSK. Using Craig's form~\eqref{eq:appA:craig} and swapping
the order of integration,
\begin{equation}
\bar P=\E_\gamma\big[\Q(\sqrt{2g\gamma})\big]=\frac1\pi\int_0^{\pi/2}M_\gamma\!\left(-\frac{g}{\sin^2\theta}\right)d\theta,
\label{eq:appA:mgfmethod}
\end{equation}
a single finite integral of the MGF of $\gamma$. Table~\ref{tab:appA:mgf} gives the MGFs; for
$L$ independent diversity branches with maximal-ratio combining, the MGF is the product of the
branch MGFs. The approach is due to Simon and Alouini, whose book is the standard reference.

\begin{table}[htbp]\centering\small
\caption{MGFs of the SNR $\gamma$ (mean $\bar\gamma$) for common fading models.}\label{tab:appA:mgf}
\renewcommand{\arraystretch}{1.5}
\begin{tabular}{@{}lll@{}}
\toprule
Fading & $M_\gamma(s)$ & BPSK average $P_b$ (closed form)\\
\midrule
Rayleigh & $\dfrac{1}{1-s\bar\gamma}$ & $\tfrac12\Big(1-\sqrt{\tfrac{\bar\gamma}{1+\bar\gamma}}\Big)\approx\dfrac1{4\bar\gamma}$\\
Nakagami-$m$ & $\Big(1-\dfrac{s\bar\gamma}{m}\Big)^{-m}$ & $\propto\bar\gamma^{-m}$ at high SNR\\
Rice ($K$) & $\dfrac{1+K}{1+K-s\bar\gamma}\exp\!\Big(\dfrac{Ks\bar\gamma}{1+K-s\bar\gamma}\Big)$ & via~\eqref{eq:appA:mgfmethod}\\
$L$-branch MRC, i.i.d.\ Rayleigh & $(1-s\bar\gamma)^{-L}$ & $\approx\binom{2L-1}{L}(4\bar\gamma)^{-L}$\\
\bottomrule
\end{tabular}
\end{table}

\begin{worked}[BPSK in Rayleigh fading at 20 dB]
At $\bar\gamma=20$\,dB $=100$, the exact Rayleigh result gives
$P_b=\tfrac12(1-\sqrt{100/101})=2.48\times10^{-3}$, and the high-SNR form $1/(4\bar\gamma)$ gives
$2.5\times10^{-3}$. The same BPSK link without fading at 20\,dB would have
$P_b=\Q(\sqrt{200})\approx10^{-45}$. The 40-odd orders of magnitude are the price of deep fades,
and they explain why every wireless system since GSM spends so much effort on diversity, coding
and interleaving. With two-branch MRC the approximation in Table~\ref{tab:appA:mgf} gives
$3/(400)^2\approx1.9\times10^{-5}$: a hundredfold improvement from one extra antenna.
\end{worked}

\subsection{Gaussian random vectors}

A real Gaussian vector $\vect x\sim\mathcal N(\bm{\mu},\vect C)$ of dimension $n$ has density
\begin{equation}
f(\vect x)=\frac{1}{(2\pi)^{n/2}\det(\vect C)^{1/2}}\exp\!\Big(-\tfrac12(\vect x-\bm{\mu})^T\vect C^{-1}(\vect x-\bm{\mu})\Big).
\end{equation}
Everything about it is determined by $\bm{\mu}$ and $\vect C$, and Gaussianity survives every
linear operation:
\begin{itemize}[leftmargin=1.5em,itemsep=2pt]
\item \emph{Linear transforms}: $\vect y=\vect A\vect x+\vect b\sim\mathcal N(\vect A\bm{\mu}+\vect b,\;\vect A\vect C\vect A^T)$.
\item \emph{Whitening}: with $\vect C=\vect L\vect L^T$ (Cholesky), $\vect L^{-1}(\vect x-\bm{\mu})\sim\mathcal N(\vect0,\vect I)$.
Generating correlated noise reverses the step: $\vect x=\bm{\mu}+\vect L\vect w$ with white $\vect w$.
\item \emph{Uncorrelated implies independent} for jointly Gaussian variables (and only for them).
\item \emph{Conditioning}: for jointly Gaussian $(\vect x,\vect y)$,
\begin{equation}
\E[\vect x|\vect y]=\bm{\mu}_x+\vect C_{xy}\vect C_{yy}^{-1}(\vect y-\bm{\mu}_y),\qquad
\operatorname{Cov}[\vect x|\vect y]=\vect C_{xx}-\vect C_{xy}\vect C_{yy}^{-1}\vect C_{yx}.
\label{eq:appA:condgauss}
\end{equation}
The conditional mean is linear in $\vect y$; it is the MMSE estimator, which is why LMMSE channel
estimation and MMSE equalisation are optimal, not merely convenient, in Gaussian settings
(Chapters~\ref{ch:equalization} and~\ref{ch:ofdm}).
\item \emph{Quadratic forms}: $\E[\vect x^T\vect A\vect x]=\operatorname{tr}(\vect A\vect C)+\bm{\mu}^T\vect A\bm{\mu}$.
\end{itemize}

\paragraph{Circularly symmetric complex Gaussian vectors.} A complex vector $\vect z$ is
\emph{proper} (circularly symmetric) if $\E[(\vect z-\bm{\mu})(\vect z-\bm{\mu})^T]=\vect 0$,
i.e.\ its pseudo-covariance vanishes; then its real and imaginary parts have equal covariance and
a skew-symmetric cross-covariance, and $\vect z\sim\mathcal{CN}(\bm{\mu},\vect C)$ has density
\begin{equation}
f(\vect z)=\frac{1}{\pi^n\det\vect C}\exp\!\big(-(\vect z-\bm{\mu})^H\vect C^{-1}(\vect z-\bm{\mu})\big).
\label{eq:appA:cn}
\end{equation}
Thermal noise at complex baseband and the entries of a Rayleigh MIMO channel are of this type.
Linear maps preserve it, $\vect A\vect z\sim\mathcal{CN}(\vect A\bm{\mu},\vect A\vect C\vect A^H)$, and
$\mathcal{CN}(\vect0,\sigma^2\vect I)$ is invariant under any unitary transformation, which is why
white noise stays white after an FFT, a rotation or the SVD of a MIMO channel. For a scalar
$z\sim\mathcal{CN}(0,\sigma^2)$, $\E|z|^4=2\sigma^4$. The differential entropy is
$\log_2\det(\pi e\vect C)$ bits, the source of the $\log\det$ form of MIMO capacity
(Chapter~\ref{ch:mimo}). Improper signals (BPSK, offset QPSK, some interference) need the
pseudo-covariance as well; \emph{widely linear} processing exploits it.

%======================================================================
\section{Random processes}
\label{sec:appA:processes}
%======================================================================

A process $x(t)$ is \term{wide-sense stationary} (WSS) if its mean is constant and its
autocorrelation $R_x(\tau)=\E[x(t+\tau)x^*(t)]$ depends only on the lag $\tau$. The
\term{Wiener--Khinchin theorem} identifies its power spectral density as
\begin{equation}
S_x(f)=\int_{-\infty}^{\infty}R_x(\tau)\,e^{-j2\pi f\tau}d\tau,\qquad
\E|x(t)|^2=R_x(0)=\int S_x(f)\,df.
\label{eq:appA:wk}
\end{equation}
For a real process $S_x$ is real, even and non-negative. Passing $x$ through an LTI filter $h$ gives
\begin{equation}
S_y(f)=|H(f)|^2S_x(f),\qquad R_{yx}(\tau)=h(\tau)*R_x(\tau),\qquad R_y(\tau)=h(\tau)*h^*(-\tau)*R_x(\tau).
\label{eq:appA:lti}
\end{equation}
The cross-correlation result is the basis of channel sounding: drive a system with white noise
(or a sequence with an impulse-like autocorrelation) and the input--output cross-correlation is
the impulse response.

\paragraph{White noise and noise bandwidth.} Real thermal noise is modelled as white with
two-sided PSD $N_0/2$ and $R(\tau)=\tfrac{N_0}2\delta(\tau)$, with $N_0=kT$
(Chapter~\ref{ch:noise}). The \term{noise-equivalent bandwidth} of a filter is the width of the
ideal rectangle with the same peak gain that passes the same noise power:
\begin{equation}
B_N=\frac{1}{|H(f_0)|^2}\int_0^\infty|H(f)|^2df .
\label{eq:appA:enbw}
\end{equation}
A one-pole RC low-pass has $B_N=\tfrac\pi2f_{3\,\mathrm{dB}}$; an $n$-pole Butterworth approaches
$f_{3\,\mathrm{dB}}$ as $n$ grows. A root-raised-cosine matched filter, considered at complex
baseband over positive and negative frequencies, passes exactly $N_0R_s$ of noise power whatever
its roll-off, which is why $E_s/N_0$ and the SNR at the sampler coincide.

\paragraph{Bandpass noise.} Gaussian noise filtered to a band around $f_c$ can be written
$n(t)=n_I(t)\cos2\pi f_ct-n_Q(t)\sin2\pi f_ct$, where $n_I$ and $n_Q$ are low-pass, have the same
variance as $n(t)$ and, for a PSD symmetric about $f_c$, are independent. Their PSD is
$S_n(f-f_c)+S_n(f+f_c)$ restricted to low frequencies: white bandpass noise of PSD $N_0/2$ and
width $B$ gives $n_I$ and $n_Q$ of PSD $N_0$ over $|f|<B/2$, each with power $N_0B$. The envelope
$\sqrt{n_I^2+n_Q^2}$ is Rayleigh and the phase uniform.

\paragraph{Discrete time.} For a WSS sequence, $R_x[m]=\E[x[n+m]x^*[n]]$ and
$S_x(e^{j\omega})=\sum_mR_x[m]e^{-j\omega m}$. A first-order autoregressive process
$x[n]=ax[n-1]+w[n]$ with white $w$ of variance $\sigma_w^2$ has $R_x[m]=\sigma_w^2a^{|m|}/(1-a^2)$
and $S_x=\sigma_w^2/|1-ae^{-j\omega}|^2$; it is the usual model for a slowly varying channel
gain tracked by a Kalman filter. Sampled white noise of PSD $N_0$ (complex) at rate $f_s$ after an
ideal anti-alias filter has variance $N_0f_s$ per sample.

\paragraph{Ergodicity.} Time averages of one realisation converge to ensemble averages for an
ergodic process. Every Monte Carlo BER simulation and every spectrum analyser assume it.

\paragraph{Phase noise.} Free-running oscillator phase is well modelled as a Wiener process:
$\phi(t+\tau)-\phi(t)\sim\mathcal N(0,2\pi\Delta\nu|\tau|)$. The resulting carrier
$e^{j\phi(t)}$ has a Lorentzian spectrum whose full width at half maximum is $\Delta\nu$, the
\emph{linewidth}. This is the model of \texttt{commlib.phase\_noise} and of laser phase noise in
coherent optics (Chapters~\ref{ch:sdr} and~\ref{ch:wireline}).

\paragraph{Cyclostationarity.} Modulated signals are not stationary: a PAM signal
$s(t)=\sum_ka_kp(t-kT)$ has an autocorrelation $R_s(t+\tau,t)$ that is periodic in $t$ with
period $T$. Such a process is \term{cyclostationary}. Averaging over one period gives the
stationary PSD used throughout the book,
\begin{equation}
\bar S_s(f)=\frac1T|P(f)|^2S_a(e^{j2\pi fT}),\qquad S_a(e^{j\omega})=\sum_mR_a[m]e^{-j\omega m},
\label{eq:appA:pampsd}
\end{equation}
for zero-mean symbols (Chapter~\ref{ch:baseband} adds the spectral lines of a non-zero mean). The
periodic part is described by the cyclic autocorrelation
$R^\alpha(\tau)=\frac1T\int_0^TR_s(t+\tau,t)e^{-j2\pi\alpha t}dt$ at cycle frequencies
$\alpha=k/T$, and its Fourier transform, the spectral correlation function, has features at the
symbol rate and twice the carrier that noise does not. Squaring a signal to recover its symbol
clock, the Gardner detector and cyclostationary feature detection in spectrum sensing all
exploit them (Chapter~\ref{ch:sync}).

%======================================================================
\section{Linear algebra for communications}
\label{sec:appA:linalg}
%======================================================================

Equalisers, MIMO detectors, precoders, channel estimators and LDPC decoders are all matrix
computations. This section lists what the book uses, with complex matrices throughout.

\subsection{Special matrices and decompositions}

A square matrix $\vect A$ is \emph{Hermitian} if $\vect A^H=\vect A$ (real symmetric in the real
case), \emph{unitary} if $\vect A^H\vect A=\vect I$, \emph{normal} if $\vect A\vect A^H=\vect A^H\vect A$,
and \emph{positive definite} ($\vect A\succ0$) if $\vect x^H\vect A\vect x>0$ for every
$\vect x\ne\vect 0$. Covariance and Gram matrices ($\vect H^H\vect H$) are Hermitian positive
semidefinite. Unitary matrices preserve lengths and inner products: they are the rotations and
reflections of $\C^n$, and the unitary DFT matrix is one of them.

\begin{description}[leftmargin=1.2em,itemsep=3pt,font=\normalfont\itshape]
\item[Eigendecomposition.] A Hermitian matrix has real eigenvalues and orthonormal eigenvectors:
$\vect A=\vect U\bm{\Lambda}\vect U^H$. Then $\operatorname{tr}\vect A=\sum\lambda_i$,
$\det\vect A=\prod\lambda_i$, and the Rayleigh quotient satisfies
$\lambda_{\min}\le\vect x^H\vect A\vect x/\vect x^H\vect x\le\lambda_{\max}$. The eigenvalue spread
$\lambda_{\max}/\lambda_{\min}$ of the input correlation matrix sets the convergence speed of LMS
(Chapter~\ref{ch:equalization}); the dominant eigenvector of a spatial covariance is the
best fixed beam (Chapter~\ref{ch:mimo}).
\item[Singular value decomposition.] Every $m\times n$ matrix factors as
$\vect H=\vect U\bm{\Sigma}\vect V^H$ with unitary $\vect U$, $\vect V$ and non-negative singular
values $\sigma_1\ge\sigma_2\ge\dots$ on the diagonal of $\bm{\Sigma}$. The rank is the number of
non-zero $\sigma_i$; the $\sigma_i^2$ are the eigenvalues of $\vect H^H\vect H$. Precoding with
$\vect V$ and combining with $\vect U^H$ turns a MIMO channel into $\operatorname{rank}\vect H$
parallel scalar channels with gains $\sigma_i$, the basis of MIMO capacity and water-filling
(Chapter~\ref{ch:mimo}). The best rank-$r$ approximation (Eckart--Young) keeps the $r$ largest
terms. The condition number $\sigma_{\max}/\sigma_{\min}$ predicts the noise enhancement of
zero-forcing.
\item[QR and Cholesky.] $\vect H=\vect Q\vect R$ with unitary $\vect Q$ and upper-triangular
$\vect R$ turns detection into back-substitution: it is the first step of successive
interference cancellation and the sphere decoder. A positive definite $\vect A$ factors as
$\vect A=\vect L\vect L^H$ (Cholesky), the cheapest way to solve normal equations and to whiten
correlated noise.
\item[Pseudo-inverse and least squares.] For full column rank $\vect A$, the least-squares
solution of $\vect y\approx\vect A\vect x$ is $\hat{\vect x}=\vect A^+\vect y$ with
$\vect A^+=(\vect A^H\vect A)^{-1}\vect A^H=\vect V\bm{\Sigma}^+\vect U^H$, and
$\vect P=\vect A\vect A^+$ is the orthogonal projection onto the column space. The regularised
(MMSE) version is $(\vect A^H\vect A+\sigma^2\vect I)^{-1}\vect A^H\vect y
=\vect A^H(\vect A\vect A^H+\sigma^2\vect I)^{-1}\vect y$, the second form being cheaper when $\vect A$
is wide.
\item[Norms.] $\|\vect A\|_F^2=\operatorname{tr}(\vect A^H\vect A)=\sum_{ij}|a_{ij}|^2=\sum\sigma_i^2$;
the spectral norm $\|\vect A\|_2=\sigma_{\max}$.
\end{description}

\subsection{Identities}

\begin{table}[htbp]\centering\small
\caption{Matrix identities. Dimensions are assumed compatible and inverses to exist.}\label{tab:appA:matid}
\renewcommand{\arraystretch}{1.45}
\begin{tabular}{@{}ll@{}}
\toprule
Identity & Where it is used\\
\midrule
$(\vect A+\vect U\vect C\vect V)^{-1}=\vect A^{-1}-\vect A^{-1}\vect U(\vect C^{-1}+\vect V\vect A^{-1}\vect U)^{-1}\vect V\vect A^{-1}$ & matrix inversion lemma (Woodbury)\\
$(\vect A+\vect u\vect v^H)^{-1}=\vect A^{-1}-\dfrac{\vect A^{-1}\vect u\vect v^H\vect A^{-1}}{1+\vect v^H\vect A^{-1}\vect u}$ & Sherman--Morrison: the RLS update\\
$(\vect I+\vect A\vect B)^{-1}\vect A=\vect A(\vect I+\vect B\vect A)^{-1}$ & push-through: MMSE forms\\
$\det(\vect I_m+\vect A\vect B)=\det(\vect I_n+\vect B\vect A)$ & Sylvester: MIMO capacity either side\\
$\operatorname{tr}(\vect A\vect B\vect C)=\operatorname{tr}(\vect C\vect A\vect B)$ & cyclic trace\\
$\begin{bmatrix}\vect A&\vect B\\\vect C&\vect D\end{bmatrix}^{-1}=\begin{bmatrix}\vect S^{-1}&-\vect S^{-1}\vect B\vect D^{-1}\\-\vect D^{-1}\vect C\vect S^{-1}&\vect D^{-1}+\vect D^{-1}\vect C\vect S^{-1}\vect B\vect D^{-1}\end{bmatrix}$ & block inverse, $\vect S=\vect A-\vect B\vect D^{-1}\vect C$\\
$\det\begin{bmatrix}\vect A&\vect B\\\vect C&\vect D\end{bmatrix}=\det(\vect D)\det(\vect A-\vect B\vect D^{-1}\vect C)$ & Schur complement\\
\midrule
$\operatorname{vec}(\vect A\vect X\vect B)=(\vect B^T\otimes\vect A)\operatorname{vec}(\vect X)$ & channel estimation, Kronecker models\\
$(\vect A\otimes\vect B)(\vect C\otimes\vect D)=\vect A\vect C\otimes\vect B\vect D$ & mixed product\\
$(\vect A\otimes\vect B)^{-1}=\vect A^{-1}\otimes\vect B^{-1}$,\; $(\vect A\otimes\vect B)^H=\vect A^H\otimes\vect B^H$ & \\
$\operatorname{tr}(\vect A\otimes\vect B)=\operatorname{tr}\vect A\operatorname{tr}\vect B$,\; $\det(\vect A_{n\times n}\otimes\vect B_{m\times m})=\det(\vect A)^m\det(\vect B)^n$ & \\
$\operatorname{tr}(\vect A^H\vect B)=\operatorname{vec}(\vect A)^H\operatorname{vec}(\vect B)$ & \\
eigenvalues of $\vect A\otimes\vect B$ are all products $\lambda_i(\vect A)\lambda_j(\vect B)$ & planar arrays, 2-D DFT\\
\bottomrule
\end{tabular}
\end{table}

\paragraph{The Kronecker channel model.} A MIMO channel with receive correlation $\vect R_r$ and
transmit correlation $\vect R_t$ is often modelled as
$\vect H=\vect R_r^{1/2}\vect G\vect R_t^{1/2}$ with $\vect G$ i.i.d.\ $\mathcal{CN}(0,1)$. By the first
vec identity, $\operatorname{Cov}[\operatorname{vec}\vect H]=\vect R_t^T\otimes\vect R_r$
(\texttt{commlib.correlated\_mimo}). The array response of an $N_x\times N_y$ planar array is
likewise a Kronecker product of two linear-array responses.

\subsection{Complex gradients (Wirtinger calculus)}

A real cost $J(\vect z)$ of a complex vector, such as a mean-square error, is not
complex-differentiable, but it can be optimised by treating $\vect z$ and $\vect z^*$ as
independent variables. With $z=x+jy$,
\begin{equation}
\frac{\partial}{\partial z}=\frac12\Big(\frac{\partial}{\partial x}-j\frac{\partial}{\partial y}\Big),\qquad
\frac{\partial}{\partial z^*}=\frac12\Big(\frac{\partial}{\partial x}+j\frac{\partial}{\partial y}\Big).
\label{eq:appA:wirtinger}
\end{equation}
A stationary point of a real $J$ satisfies $\partial J/\partial\vect z^*=\vect0$, and the direction of
steepest \emph{ascent} is $\partial J/\partial\vect z^*$ (not $\partial J/\partial\vect z$).
Table~\ref{tab:appA:wirtinger} gives the derivatives that recur, in the convention that the gradient
with respect to a column vector is a column vector.

\begin{table}[htbp]\centering\small
\caption{Wirtinger and matrix derivatives. $\vect a$, $\vect A$ constant.}\label{tab:appA:wirtinger}
\renewcommand{\arraystretch}{1.35}
\begin{tabular}{@{}lll@{}}
\toprule
$J$ & $\partial J/\partial\vect z^*$ & Notes\\
\midrule
$\vect a^H\vect z$ & $\vect 0$ & depends on $\vect z$ only\\
$\vect z^H\vect a$ & $\vect a$ & \\
$\vect z^H\vect A\vect z$ & $\vect A\vect z$ & $\vect A$ Hermitian: real quadratic form\\
$\|\vect y-\vect H\vect z\|^2$ & $-\vect H^H(\vect y-\vect H\vect z)$ & zero gives the normal equations\\
$\E|d-\vect w^H\vect u|^2$ (in $\vect w$) & $\vect R\vect w-\vect p$ & $\vect R=\E\vect u\vect u^H$, $\vect p=\E\vect u d^*$; Wiener: $\vect w=\vect R^{-1}\vect p$\\
\midrule
$J(\vect X)$ & $\partial J/\partial\vect X$ & \\
\midrule
$\operatorname{tr}(\vect A\vect X)$ & $\vect A^T$ & \\
$\ln\det\vect X$ & $\vect X^{-T}$ & log-likelihoods, capacity\\
$\operatorname{tr}(\vect X^H\vect A\vect X)$ (in $\vect X^*$) & $\vect A\vect X$ & \\
$\vect x^T\vect A\vect x$ (real) & $(\vect A+\vect A^T)\vect x$ & \\
\bottomrule
\end{tabular}
\end{table}

\begin{worked}[Deriving LMS in two lines]
An adaptive equaliser with taps $\vect w$ and input vector $\vect u_n$ produces $\vect w^H\vect u_n$
and makes the error $e_n=d_n-\vect w^H\vect u_n$. The instantaneous cost is $|e_n|^2=e_ne_n^*$, and
since $e_n$ depends on $\vect w^*$ through $\vect w^H\vect u_n=\sum w_i^*u_i$ while $e_n^*$ depends
only on $\vect w$, $\partial|e_n|^2/\partial\vect w^*=-\vect u_ne_n^*$. A step down the gradient is
\[
\vect w\leftarrow\vect w+\mu\,\vect u_n\,e_n^*,
\]
the LMS algorithm of Chapter~\ref{ch:equalization} (and of \texttt{commlib.nlms\_equalizer}, which
divides $\mu$ by $\|\vect u_n\|^2$). Using the expected cost instead gives $\vect R\vect w=\vect p$,
the Wiener--Hopf equations.
\end{worked}

\subsection{Toeplitz and circulant matrices}
\label{sec:appA:circulant}

Linear convolution of a length-$N$ input with a length-$L$ channel is multiplication by an
$(N+L-1)\times N$ \term{Toeplitz} matrix (constant along diagonals): this is
\texttt{commlib.conv\_matrix}, and the autocorrelation matrix of any WSS sequence is Hermitian
Toeplitz. Toeplitz systems can be solved in $O(N^2)$ operations by the Levinson--Durbin recursion
(the engine of LPC speech coding, Chapter~\ref{ch:sourcecoding}), and by Szeg\H{o}'s theorem the
eigenvalues of a large Toeplitz covariance matrix are distributed like samples of the process's
PSD.

A \term{circulant} matrix is Toeplitz with wrap-around: each row is the previous one rotated right,
so $\vect C\vect x$ is the circular convolution of $\vect x$ with the first column $\vect c$. Every
circulant matrix is diagonalised by the DFT:
\begin{equation}
\vect C=\vect F^H\operatorname{diag}(\lambda_0,\dots,\lambda_{N-1})\,\vect F,\qquad
\lambda_k=\sum_{n=0}^{N-1}c_n\,e^{-j2\pi kn/N},\qquad [\vect F]_{kn}=\tfrac1{\sqrt N}e^{-j2\pi kn/N}.
\label{eq:appA:circulant}
\end{equation}
The eigenvectors are the complex exponentials, independent of $\vect c$, and the eigenvalues are the
channel's frequency response at the DFT bins. Adding a cyclic prefix makes the channel matrix seen
by the receiver circulant, so the receiver's FFT diagonalises it: OFDM in one line
(Chapter~\ref{ch:ofdm}). The same structure lets single-carrier frequency-domain equalisers and
DFT-spread OFDM invert a channel with one division per bin, and makes the 2-D DFT diagonalise the
doubly-circulant matrices of OTFS.

%======================================================================
\section{Optimisation}
\label{sec:appA:optimisation}
%======================================================================

\paragraph{Unconstrained.} A differentiable function has zero gradient at an interior optimum;
for a convex function (non-negative-definite Hessian everywhere) every such point is a global
minimum. Gradient descent $\vect x\leftarrow\vect x-\mu\nabla J$ converges for small enough $\mu$ on
smooth convex problems; Newton's method $\vect x\leftarrow\vect x-\nabla^2J^{-1}\nabla J$ converges
quadratically near the optimum. LMS is stochastic gradient descent and RLS is a recursive
Newton method on the least-squares cost.

\paragraph{Equality constraints: Lagrange multipliers.} To minimise $J(\vect x)$ subject to
$g_i(\vect x)=0$, form the Lagrangian $\mathcal L=J+\sum_i\nu_ig_i$ and set its gradient to zero. The
multiplier $\nu_i$ measures how fast the optimum changes as constraint $i$ is relaxed.

\begin{worked}[The minimum-variance (MVDR) beamformer]
Minimise the output power $\vect w^H\vect R\vect w$ of an array subject to unit gain in the look
direction, $\vect w^H\vect a=1$. The Lagrangian
$\vect w^H\vect R\vect w+\nu(1-\vect w^H\vect a)+\nu^*(1-\vect a^H\vect w)$ has
$\partial/\partial\vect w^*=\vect R\vect w-\nu\vect a=\vect0$, so $\vect w=\nu\vect R^{-1}\vect a$.
The constraint fixes $\nu=1/(\vect a^H\vect R^{-1}\vect a)$:
\[
\vect w_{\mathrm{MVDR}}=\frac{\vect R^{-1}\vect a}{\vect a^H\vect R^{-1}\vect a}.
\]
This is the Capon beamformer of Chapter~\ref{ch:mimo}; with white noise ($\vect R\propto\vect I$)
it reduces to the matched (conventional) beamformer $\vect a/\|\vect a\|^2$.
\end{worked}

\paragraph{Inequality constraints: the KKT conditions.} For the convex problem
$\min J(\vect x)$ subject to $h_i(\vect x)\le0$ and $g_j(\vect x)=0$ (convex $J$ and $h_i$, affine
$g_j$, and a strictly feasible point), $\vect x^\star$ is optimal if and only if there are multipliers
$\lambda_i$ and $\nu_j$ with
\begin{equation}
\nabla J(\vect x^\star)+\sum_i\lambda_i\nabla h_i(\vect x^\star)+\sum_j\nu_j\nabla g_j(\vect x^\star)=\vect0,\quad
h_i\le0,\; g_j=0,\quad \lambda_i\ge0,\quad \lambda_ih_i(\vect x^\star)=0.
\label{eq:appA:kkt}
\end{equation}
These are the Karush--Kuhn--Tucker conditions: stationarity, primal feasibility, dual feasibility and
\emph{complementary slackness}. The last one says each inequality is either active or has a zero
multiplier.

\paragraph{Water-filling.} Allocate total power $P$ over $K$ parallel channels with gains $g_k$
(SNR per unit power) to maximise $\sum_k\log_2(1+g_kp_k)$ subject to $\sum p_k\le P$ and
$p_k\ge0$. The problem is concave; applying~\eqref{eq:appA:kkt} with multiplier $\nu$ for the
power constraint and $\lambda_k$ for $p_k\ge0$ gives $g_k/(1+g_kp_k)=(\nu-\lambda_k)\ln 2$, and
complementary slackness turns this into
\begin{equation}
p_k=\Big(\mu-\frac{1}{g_k}\Big)^+,\qquad\text{with the ``water level'' }\mu\text{ chosen so that }\sum_kp_k=P.
\label{eq:appA:waterfill}
\end{equation}
Picture the inverse gains $1/g_k$ as the uneven floor of a vessel and pour in $P$ units of water:
deep (good) channels get more, and channels whose floor lies above the water get nothing.
It is the capacity-achieving allocation over frequency for a known ISI channel and over
eigenmodes for MIMO (Chapters~\ref{ch:infotheory} and~\ref{ch:mimo}); DSL bit loading is its
discrete cousin (Chapter~\ref{ch:wireline}).

\begin{worked}[Water-filling over three channels]
Let $g=(10,5,1)$ and $P=1$, so the floor is $1/g=(0.1,0.2,1)$. Try all three channels: the water
level would be $\mu=(1+0.1+0.2+1)/3=0.767$, below the third floor, so drop channel 3 and retry
with two: $\mu=(1+0.1+0.2)/2=0.65$, which is above both remaining floors. The allocation is
$p=(0.55,\,0.45,\,0)$ and the rate $\log_2(6.5)+\log_2(3.25)=4.40$\,b/s/Hz, against 3.95 with
equal power $1/3$ per channel. \texttt{commlib.infotheory.waterfill([0.1, 0.2, 1.0], 1.0)}
returns the same allocation and level.
\end{worked}

Other optimisation patterns in the book: max--min and proportional fairness in scheduling
(maximise $\sum\log r_k$, Chapter~\ref{ch:multipleaccess}); the Blahut--Arimoto alternating
maximisation for channel capacity (Chapter~\ref{ch:infotheory}); and the Cauchy--Schwarz
inequality $|\langle\vect a,\vect b\rangle|\le\|\vect a\|\|\vect b\|$, with equality iff
$\vect a\propto\vect b$, which solves the matched-filter, maximal-ratio-combining and
beam-steering problems at a stroke.

%======================================================================
\section{Decibels and the units engineers mix}
\label{sec:appA:db}
%======================================================================

A power ratio in decibels is $10\log_{10}(P_2/P_1)$; for voltages or fields across the same
impedance it is $20\log_{10}(V_2/V_1)$. Adding a suffix makes the decibel an \emph{absolute}
level by naming the reference. The rules that prevent most link-budget errors are:
\begin{enumerate}[leftmargin=1.5em,itemsep=2pt]
\item Gains and losses (plain dB) add to levels: $\mathrm{dBm}+\mathrm{dB}=\mathrm{dBm}$.
\item Two absolute levels never add in dB. Two 0\,dBm signals together make 3\,dBm, not
0\,dBm\,+\,0\,dBm; convert to watts, add, and convert back.
\item Subtracting levels with the same reference gives a ratio in dB: signal (dBm) minus noise
(dBm) is SNR (dB).
\item Densities carry their bandwidth: a level in dBm/Hz plus $10\log_{10}B$ (dB-Hz) is a level in dBm.
\end{enumerate}

\begin{table}[htbp]\centering\small
\caption{Decibel units.}\label{tab:appA:dbunits}
\begin{tabular}{@{}lp{0.38\linewidth}p{0.36\linewidth}@{}}
\toprule
Unit & Reference & Typical use\\
\midrule
dB & a ratio (no reference) & gain, loss, SNR, $E_b/N_0$\\
dBm, dBW & 1\,mW; 1\,W ($\mathrm{dBW}=\mathrm{dBm}-30$) & transmit and received power\\
dBm/Hz, dBW/Hz & 1\,mW (1\,W) in 1\,Hz & noise and signal densities, $N_0$\\
dB$\mu$V & 1\,$\mu$V (in $50\,\Omega$: $\mathrm{dB}\mu\mathrm{V}=\mathrm{dBm}+107$) & receiver sensitivity, EMC\\
dB$\mu$V/m & field strength of 1\,$\mu$V/m & broadcast coverage, EMC limits\\
dBW/m$^2$ & power flux density of 1\,W/m$^2$ & satellite PFD limits\\
dBi & gain of an isotropic antenna & antenna gain, EIRP\\
dBd & gain of a half-wave dipole ($\mathrm{dBi}=\mathrm{dBd}+2.15$) & land-mobile antennas, ERP\\
dBc & the carrier (or wanted signal) power & spurs, harmonics, ACLR\\
dBc/Hz & carrier power, per Hz of offset bandwidth & oscillator phase noise\\
dB-Hz & 1\,Hz & $C/N_0$, bandwidth as $10\log_{10}B$\\
dBK & 1\,K & noise temperature, $T_{\mathrm{sys}}$\\
dB/K & ratio of gain to 1\,K & receive figure of merit $G/T$\\
dBW/K/Hz & 1\,W per kelvin per hertz & Boltzmann's constant, $-228.6$\\
dBFS & full scale of a converter & ADC/DAC levels, digital headroom\\
dBsm & 1\,m$^2$ & radar cross-section\\
dBm0, dBr & level relative to a reference point & telephone transmission plans\\
neper (Np) & natural-log ratio: $1\,\mathrm{Np}=8.686$\,dB & transmission-line attenuation\\
\bottomrule
\end{tabular}
\end{table}

\begin{table}[htbp]\centering\small
\caption{Quick conversions.}\label{tab:appA:dbquick}
\begin{tabular}{@{}rr@{\qquad}rr@{\qquad}rr@{}}
\toprule
dB & power ratio & dBm & power & dBm in 50\,$\Omega$ & RMS voltage\\
\midrule
0 & 1.000 & $+60$ & 1\,kW & $+30$ & 7.07\,V\\
1 & 1.259 & $+30$ & 1\,W & $+10$ & 707\,mV\\
2 & 1.585 & $+20$ & 100\,mW & $0$ & 224\,mV\\
3 & 1.995 & $+10$ & 10\,mW & $-10$ & 70.7\,mV\\
4 & 2.512 & $0$ & 1\,mW & $-47$ & 1\,mV\\
5 & 3.162 & $-30$ & 1\,$\mu$W & $-60$ & 224\,$\mu$V\\
6 & 3.981 & $-60$ & 1\,nW & $-107$ & 1\,$\mu$V\\
7 & 5.012 & $-90$ & 1\,pW & $-120$ & 224\,nV\\
8 & 6.310 & $-120$ & 1\,fW & & \\
9 & 7.943 & $-174$ & $kT_0$ in 1\,Hz & & \\
10 & 10.00 & & & & \\
\bottomrule
\end{tabular}
\end{table}

Mental arithmetic: 3\,dB is a factor of 2 (1.995), 10\,dB a factor of 10, 1\,dB is 26\%, 0.1\,dB
is 2.3\%, and $10\log_{10}(1+x)\approx4.34x$ for small $x$. A 6\,dB voltage step is a factor of 2
in voltage and 4 in power.

\paragraph{Noise.} The thermal noise density at the standard reference temperature
$T_0=290$\,K is
\begin{equation}
kT_0=1.380649\times10^{-23}\times290=4.00\times10^{-21}\,\mathrm{W/Hz}\;=\;-204.0\,\mathrm{dBW/Hz}\;=\;-174.0\,\mathrm{dBm/Hz},
\label{eq:appA:kt0}
\end{equation}
so a receiver of noise figure NF and bandwidth $B$ has noise floor
$-174+10\log_{10}B+\mathrm{NF}$ dBm. Table~\ref{tab:appA:noisefloor} gives the thermal floor for
common channel widths and the conversion between noise figure and noise temperature,
$T_e=T_0(F-1)$.

\begin{table}[htbp]\centering\small
\caption{Thermal noise floors ($kT_0B$) and noise figure versus noise temperature.}\label{tab:appA:noisefloor}
\begin{tabular}{@{}lr@{\qquad\qquad}rr@{}}
\toprule
Bandwidth & $kT_0B$ (dBm) & NF (dB) & $T_e$ (K)\\
\midrule
1\,Hz & $-174.0$ & 0.5 & 35\\
1\,kHz & $-144.0$ & 1.0 & 75\\
25\,kHz (narrowband FM) & $-130.0$ & 1.5 & 120\\
200\,kHz (GSM) & $-121.0$ & 2.0 & 170\\
1\,MHz & $-114.0$ & 3.0 & 289\\
3.84\,MHz (WCDMA) & $-108.2$ & 4.0 & 438\\
20\,MHz (LTE, Wi-Fi) & $-101.0$ & 6.0 & 865\\
100\,MHz (5G NR FR1) & $-94.0$ & 10.0 & 2610\\
400\,MHz (5G NR FR2) & $-88.0$ & & \\
\bottomrule
\end{tabular}
\end{table}

\paragraph{The link equation in decibels.}
\begin{gather}
P_r\,[\mathrm{dBm}]=P_t+G_t+G_r-L_{\mathrm{fs}}-L_{\mathrm{other}},\\
L_{\mathrm{fs}}=20\log_{10}d_{\mathrm{km}}+20\log_{10}f_{\mathrm{MHz}}+32.45,\\
\frac{C}{N_0}\,[\mathrm{dB\text{-}Hz}]=\mathrm{EIRP}\,[\mathrm{dBW}]-L+\frac GT\,[\mathrm{dB/K}]+228.6,\\
\frac{E_b}{N_0}=\frac{C}{N_0}-10\log_{10}R_b,\qquad
\frac{E_s}{N_0}=\frac{E_b}{N_0}+10\log_{10}(r\log_2M),
\end{gather}
with code rate $r$ and constellation size $M$ (Chapters~\ref{ch:noise}, \ref{ch:channels}
and~\ref{ch:satellite}). The received power from a field strength $E$ with an isotropic antenna is
$P_r[\mathrm{dBm}]=E[\mathrm{dB}\mu\mathrm{V/m}]-20\log_{10}f_{\mathrm{MHz}}-77.2$. Noise figures
of cascaded stages combine by Friis's formula $F=F_1+(F_2-1)/G_1+(F_3-1)/(G_1G_2)+\cdots$.

\paragraph{Mismatch.} A load with reflection coefficient $\Gamma$ has return loss
$-20\log_{10}|\Gamma|$, VSWR $(1+|\Gamma|)/(1-|\Gamma|)$ and mismatch loss
$-10\log_{10}(1-|\Gamma|^2)$ (Table~\ref{tab:appA:vswr}). A VSWR of 2 reflects 11\% of the power
but costs only half a decibel, which is why antenna specifications tolerate it.

\begin{table}[htbp]\centering\small
\caption{VSWR, return loss and mismatch loss.}\label{tab:appA:vswr}
\begin{tabular}{@{}ccccc@{}}
\toprule
VSWR & $|\Gamma|$ & Return loss (dB) & Reflected power & Mismatch loss (dB)\\
\midrule
1.1 & 0.048 & 26.4 & 0.2\% & 0.01\\
1.2 & 0.091 & 20.8 & 0.8\% & 0.04\\
1.5 & 0.200 & 14.0 & 4.0\% & 0.18\\
2.0 & 0.333 & 9.5 & 11.1\% & 0.51\\
3.0 & 0.500 & 6.0 & 25.0\% & 1.25\\
5.0 & 0.667 & 3.5 & 44.4\% & 2.55\\
\bottomrule
\end{tabular}
\end{table}

\paragraph{Converters.} dBFS is measured relative to the largest signal a converter can
represent. By the usual convention a full-scale sine is 0\,dBFS, and an ideal $N$-bit quantiser
then has $\mathrm{SNR}=6.02N+1.76$\,dB (Chapter~\ref{ch:pcm}). Some instruments reference a
full-scale \emph{square} wave instead, which moves the numbers by 3\,dB: always check.

\begin{pitfall}[Bandwidths, sides and the factor of two]
Many ``3 dB errors'' in link budgets come from mixing one-sided and two-sided densities
($N_0$ versus $N_0/2$), real and complex noise, or peak and RMS amplitudes. Settle them before
starting: in this book $N_0$ is the one-sided density, complex baseband noise has variance $N_0$
per sample at one sample per symbol, and every power is an average unless marked ``peak''.
Similarly, ``EIRP'' is referenced to an isotropic radiator, ``ERP'' to a dipole: they differ by
2.15\,dB.
\end{pitfall}

%======================================================================
\section{Constants and frequency bands}
\label{sec:appA:constants}
%======================================================================

\begin{table}[htbp]\centering\small
\caption{Physical and geophysical constants. SI defining constants are exact; Earth values are
from WGS-84 and IERS conventions.}\label{tab:appA:constants}
\footnotesize
\begin{tabular}{@{}lll@{}}
\toprule
Quantity & Value & Notes\\
\midrule
Speed of light $c$ & 299\,792\,458\,m/s (exact) & $\approx3\times10^8$; 1\,ns $\approx$ 30\,cm\\
Boltzmann constant $k$ & $1.380\,649\times10^{-23}$\,J/K (exact) & $-228.6$\,dBW/K/Hz\\
Planck constant $h$ & $6.626\,070\,15\times10^{-34}$\,J\,s (exact) & photon energy $hf$\\
Elementary charge $q$ & $1.602\,176\,634\times10^{-19}$\,C (exact) & shot noise $2qI$\\
Reference temperature $T_0$ & 290\,K & IEEE noise-figure reference\\
$kT_0$ & $4.00\times10^{-21}$\,W/Hz & $-174$\,dBm/Hz\\
Permittivity of free space $\varepsilon_0$ & $8.854\times10^{-12}$\,F/m & \\
Permeability of free space $\mu_0$ & $1.2566\times10^{-6}$\,H/m & $\approx4\pi\times10^{-7}$\\
Impedance of free space $\eta_0$ & 376.73\,$\Omega$ & $\approx120\pi$; $S=E^2/\eta_0$\\
Cosmic microwave background & 2.725\,K & cold-sky floor\\
\midrule
Earth equatorial radius (WGS-84) & 6\,378.137\,km & flattening $1/298.257$\\
Earth mean radius & 6\,371.0\,km & \\
Effective Earth radius (radio) & $\approx8\,500$\,km & $k=4/3$ standard refraction\\
Earth gravitational parameter $GM$ & $3.986\,004\,418\times10^{14}$\,m$^3$/s$^2$ & Kepler: $T=2\pi\sqrt{a^3/GM}$\\
Earth rotation rate & $7.292\,115\times10^{-5}$\,rad/s & \\
Sidereal day & 86\,164.09\,s & 23\,h\,56\,min\,4\,s\\
Geostationary orbit radius & 42\,164\,km & altitude 35\,786\,km, speed 3.075\,km/s\\
GEO one-way delay (sub-satellite point) & 119\,ms & 240--280\,ms ground--satellite--ground\\
Mean Earth--Moon distance & 384\,400\,km & \\
Astronomical unit & $1.495\,978\,707\times10^{11}$\,m (exact) & \\
\midrule
Light in silica fibre & $\approx2.04\times10^8$\,m/s & $\approx$ 4.9\,$\mu$s/km\\
GPS L1, L2, L5 carriers & 1575.42, 1227.60, 1176.45\,MHz & $154f_0$, $120f_0$, $115f_0$, $f_0=10.23$\,MHz\\
\bottomrule
\end{tabular}
\end{table}

The radio spectrum is named in two overlapping ways (Figure~\ref{fig:appA_bands}). The ITU
divides it into decades, band $N$ running from $0.3\times10^N$ to $3\times10^N$\,Hz
(Recommendation ITU-R V.431; the symbols VLF to THF belong to bands 4 to 12, and ELF, SLF and ULF
are customary names for the bands below). Radar and satellite engineers use the IEEE letter bands
of IEEE Std 521, which grew out of wartime secrecy codes. Tables~\ref{tab:appA:itubands}
and~\ref{tab:appA:ieeebands} list both with typical services.

\mdcfig{appA_bands}{The two naming systems for the radio spectrum. Top: the ITU decade bands
on a logarithmic axis from 3\,Hz to 3\,THz, with representative services. Bottom: the microwave
region from 1 to 300\,GHz with the IEEE radar letter bands.}

\begin{table}[htbp]\centering\small
\caption{ITU frequency bands (ITU-R V.431).}\label{tab:appA:itubands}
\begin{tabular}{@{}clllp{0.38\linewidth}@{}}
\toprule
Band & Symbol & Frequency & Wavelength & Typical services\\
\midrule
4 & VLF & 3--30\,kHz & 100--10\,km & submarine communication, lightning sensing (historically Omega navigation)\\
5 & LF & 30--300\,kHz & 10--1\,km & longwave broadcast, time signals (60, 77.5\,kHz), Loran-C/eLoran (100\,kHz)\\
6 & MF & 0.3--3\,MHz & 1\,km--100\,m & AM broadcast, maritime, NDBs\\
7 & HF & 3--30\,MHz & 100--10\,m & shortwave broadcast, amateur, aviation (oceanic), over-the-horizon radar\\
8 & VHF & 30--300\,MHz & 10--1\,m & FM and DAB broadcast, aviation voice, marine, TV\\
9 & UHF & 0.3--3\,GHz & 1\,m--10\,cm & TV, cellular, GNSS, Wi-Fi 2.4\,GHz, Bluetooth\\
10 & SHF & 3--30\,GHz & 10--1\,cm & 5G mid-band, Wi-Fi 5/6\,GHz, satellite, radar, microwave links\\
11 & EHF & 30--300\,GHz & 10--1\,mm & 5G FR2, E-band backhaul, automotive radar, radio astronomy\\
12 & THF & 0.3--3\,THz & 1--0.1\,mm & research: sub-THz 6G, imaging\\
\bottomrule
\end{tabular}
\end{table}

\begin{table}[htbp]\centering\small
\caption{IEEE radar letter bands (IEEE Std 521) and optical bands of single-mode fibre.}\label{tab:appA:ieeebands}
\begin{tabular}{@{}llp{0.52\linewidth}@{}}
\toprule
Band & Range & Typical use\\
\midrule
L & 1--2\,GHz & GNSS, Iridium and Inmarsat, air-traffic radar, LTE/NR bands\\
S & 2--4\,GHz & Wi-Fi 2.4\,GHz, weather and naval radar, 5G n77/n78, satellite radio\\
C & 4--8\,GHz & satellite TV/feeder (3.7--4.2/5.9--6.4\,GHz), Wi-Fi 5/6\,GHz, weather radar\\
X & 8--12\,GHz & military satellite, deep-space links, marine and police radar\\
Ku & 12--18\,GHz & direct-to-home TV, VSAT, Starlink user links\\
K & 18--27\,GHz & water-vapour absorption near 22\,GHz; Ka-band downlinks begin near 17.7\,GHz\\
Ka & 27--40\,GHz & HTS satellites, 5G FR2 bands n257, n260, n261, Ka-band uplinks\\
V & 40--75\,GHz & 60\,GHz unlicensed (oxygen absorption), inter-satellite links\\
W & 75--110\,GHz & automotive radar (76--81\,GHz), E-band backhaul (71--76/81--86\,GHz, straddling V and W)\\
mm & 110--300\,GHz & D band (110--170\,GHz) research for 6G, imaging\\
\midrule
O, E, S & 1260--1360, 1360--1460, 1460--1530\,nm & original, extended and short-wavelength fibre bands\\
C, L, U & 1530--1565, 1565--1625, 1625--1675\,nm & conventional (EDFA) band, long, ultra-long\\
\bottomrule
\end{tabular}
\end{table}

\begin{pitfall}[Which ``Ka-band''?]
Letter bands are conventions, not allocations, and different communities draw them differently:
satellite operators routinely call 17.7--21.2\,GHz ``Ka-band downlink'' though IEEE Std 521 places
it in K band, and the NATO/EU band letters (A to M) are a completely different system in which
``C band'' means 0.5--1\,GHz. Always quote the frequency in a specification; quote the letter only
as a nickname.
\end{pitfall}

%======================================================================
\section{Series, sums and integrals}
\label{sec:appA:series}
%======================================================================

\paragraph{Finite and infinite sums.}
\begin{gather}
\sum_{n=0}^{N-1}r^n=\frac{1-r^N}{1-r},\qquad \sum_{n=0}^{\infty}r^n=\frac1{1-r},\qquad
\sum_{n=0}^{\infty}nr^n=\frac{r}{(1-r)^2}\quad(|r|<1),\\
\sum_{n=0}^{N-1}e^{j\omega n}=e^{j\omega(N-1)/2}\frac{\sin(N\omega/2)}{\sin(\omega/2)},\qquad
\sum_{n=0}^{N-1}e^{j2\pi kn/N}=\begin{cases}N,&k\equiv0\pmod N\\0,&\text{otherwise,}\end{cases}\\
\sum_{n=1}^{N}n=\frac{N(N+1)}{2},\qquad\sum_{n=1}^{N}n^2=\frac{N(N+1)(2N+1)}{6},\qquad
\sum_{k=0}^{n}\binom nk a^kb^{n-k}=(a+b)^n.
\end{gather}
The second line is the orthogonality of DFT basis vectors and of OFDM subcarriers. The geometric
sum with $r=\rho$ gives the mean queue length of an M/M/1 system and the impulse response energy
of a one-pole filter.

\paragraph{Expansions for small $x$.}
\begin{gather}
e^x=1+x+\tfrac{x^2}{2}+\cdots,\qquad \ln(1+x)=x-\tfrac{x^2}2+\cdots,\qquad (1+x)^a=1+ax+\cdots,\\
\sin x=x-\tfrac{x^3}{6}+\cdots,\qquad \cos x=1-\tfrac{x^2}2+\cdots,\\
\tan^{-1}x=x-\tfrac{x^3}3+\cdots,\qquad \log_2(1+x)\approx1.443x .
\end{gather}
The last is why capacity grows linearly with power at low SNR (Chapter~\ref{ch:infotheory}).

\paragraph{Trigonometric identities.} With $e^{j\theta}=\cos\theta+j\sin\theta$:
\begin{gather}
\cos a\cos b=\tfrac12[\cos(a-b)+\cos(a+b)],\qquad \sin a\sin b=\tfrac12[\cos(a-b)-\cos(a+b)],\\
\sin a\cos b=\tfrac12[\sin(a+b)+\sin(a-b)],\qquad \cos^2a=\tfrac12(1+\cos2a),\qquad \sin^2a=\tfrac12(1-\cos2a),\\
\cos(a\pm b)=\cos a\cos b\mp\sin a\sin b,\qquad \sin(a\pm b)=\sin a\cos b\pm\cos a\sin b,\\
A\cos\theta+B\sin\theta=\sqrt{A^2+B^2}\,\cos(\theta-\tan^{-1}(B/A)).
\end{gather}
The product formulas are the whole theory of mixers: multiplying by a local oscillator produces
sum and difference frequencies (Chapters~\ref{ch:analog} and~\ref{ch:sdr}).

\paragraph{Gaussian integrals.} For $a>0$,
\begin{gather}
\int_{-\infty}^{\infty}e^{-ax^2}dx=\sqrt{\frac\pi a},\qquad
\int_{-\infty}^{\infty}x^2e^{-ax^2}dx=\frac{\sqrt\pi}{2a^{3/2}},\qquad
\int_0^\infty xe^{-ax^2}dx=\frac1{2a},\\
\int_{-\infty}^{\infty}e^{-ax^2+bx}dx=\sqrt{\frac\pi a}\,e^{b^2/4a},\\
\int_{\R^n}e^{-\frac12\vect x^T\vect A\vect x+\vect b^T\vect x}d\vect x=\frac{(2\pi)^{n/2}}{\sqrt{\det\vect A}}\,e^{\frac12\vect b^T\vect A^{-1}\vect b},\\
\int_{\C^n}e^{-\vect z^H\vect A\vect z}d\vect z=\frac{\pi^n}{\det\vect A}.
\end{gather}
Completing the square in the exponent, the trick behind all of these, is also how the Gaussian
pdf's MGF and the conditional-Gaussian result~\eqref{eq:appA:condgauss} are derived.

\paragraph{Other recurring integrals.}
\begin{gather}
\int_{-\infty}^{\infty}\sinc x\,dx=\int_{-\infty}^{\infty}\sinc^2x\,dx=1,\qquad
\int_0^\infty x^ne^{-ax}dx=\frac{n!}{a^{n+1}},\qquad
\int_0^\infty e^{-ax}\cos bx\,dx=\frac{a}{a^2+b^2},\\
\int_0^{2\pi}e^{x\cos\theta}d\theta=2\pi I_0(x),\qquad
\int_0^\infty xe^{-ax^2}I_0(bx)\,dx=\frac{1}{2a}e^{b^2/4a},\qquad
\int_{-\infty}^{\infty}\frac{dx}{1+x^2}=\pi.
\end{gather}
The fifth integral shows that the Rice density integrates to one, and is the key step in the
error rate of noncoherent FSK.

\paragraph{Inequalities.} $\ln x\le x-1$ (with equality at $x=1$; it proves that relative entropy
is non-negative); Cauchy--Schwarz; the triangle inequality; Jensen's inequality; and
Hadamard's inequality $\det\vect A\le\prod_ia_{ii}$ for positive semidefinite $\vect A$, which
shows that a $\log\det$ capacity is largest when the transmit covariance is aligned with the
channel's eigenvectors so that the streams do not interfere (Chapter~\ref{ch:mimo}). For binomial
coefficients, $\binom nk\le2^{nh(k/n)}$ with $h$ the binary entropy function, the counting
argument behind the channel coding theorem (Chapter~\ref{ch:infotheory}).

\furtherreading
\begin{itemize}
\item F.~W.~J.~Olver et al.\ (eds.), \emph{NIST Digital Library of Mathematical Functions},
\url{https://dlmf.nist.gov}. The modern, free successor to Abramowitz and Stegun's
\emph{Handbook of Mathematical Functions} (1964): definitions, series, asymptotics and identities
for every special function in this appendix.
\item I.~S.~Gradshteyn and I.~M.~Ryzhik, \emph{Table of Integrals, Series, and Products}, 8th ed.,
Academic Press, 2014. When an integral will not yield, it is probably in here.
\item A.~Papoulis and S.~U.~Pillai, \emph{Probability, Random Variables and Stochastic Processes},
4th ed., McGraw-Hill, 2002. The engineering standard for distributions, transformations and random
processes, including cyclostationarity.
\item M.~K.~Simon and M.-S.~Alouini, \emph{Digital Communication over Fading Channels}, 2nd ed.,
Wiley, 2005. Craig's form, the MGF method, Marcum Q bounds and closed-form error rates for every
fading model.
\item J.~W.~Craig, ``A new, simple and exact result for calculating the probability of error for
two-dimensional signal constellations,'' \emph{Proc.\ IEEE MILCOM}, 1991; and M.~Chiani, D.~Dardari
and M.~K.~Simon, ``New exponential bounds and approximations for the computation of error
probability in fading channels,'' \emph{IEEE Trans.\ Wireless Communications}, vol.~2, no.~4, 2003.
\item R.~A.~Horn and C.~R.~Johnson, \emph{Matrix Analysis}, 2nd ed., Cambridge University Press, 2013;
and K.~B.~Petersen and M.~S.~Pedersen, \emph{The Matrix Cookbook} (freely available), a compact list
of identities and derivatives.
\item A.~Hj{\o}rungnes, \emph{Complex-Valued Matrix Derivatives}, Cambridge University Press, 2011.
Wirtinger calculus done carefully, with applications to signal processing and communications.
\item S.~Boyd and L.~Vandenberghe, \emph{Convex Optimization}, Cambridge University Press, 2004 (free
online). Duality, KKT conditions and water-filling as a worked example.
\item R.~M.~Gray, ``Toeplitz and circulant matrices: a review,'' \emph{Foundations and Trends in
Communications and Information Theory}, vol.~2, no.~3, 2006. Szeg\H{o}'s theorem and why circulants
approximate Toeplitz matrices.
\item F.~J.~Harris, ``On the use of windows for harmonic analysis with the discrete Fourier
transform,'' \emph{Proc.\ IEEE}, vol.~66, no.~1, 1978. The source of Table~\ref{tab:appA:windows}.
\item R.~N.~Bracewell, \emph{The Fourier Transform and Its Applications}, 3rd ed., McGraw-Hill, 2000.
A physical, picture-rich treatment of transforms, pairs and sampling.
\end{itemize}
